\chapter{Thermal Effects on Breakdown and Transition}\label{chapter:breakdown}

The development of turbulence in hypersonic boundary layers arises through a complex sequence of processes, beginning with the linear growth of instability modes and culminating in nonlinear breakdown. Whilst the second Mack mode is well captured by \gls{LST} during its early amplification, the subsequent nonlinear dynamics that drive complete transition remain difficult to predict with comparable accuracy. The breakdown of coherent instability structures into turbulence carries far-reaching engineering consequences, determining the magnitude of heating, frictional drag, and the aerothermal performance envelope of hypersonic vehicles. A defining feature of this regime is the overshoot in surface quantities such as the skin-friction coefficient and Stanton number. Rather than the gradual increase predicted by linear theory, both have been observed experimentally and numerically to exceed their fully turbulent levels during transition before relaxing downstream \citep{Holden1972}. This behaviour stems from the abrupt, localised enhancement of mixing and momentum transport during vortex breakdown, which injects turbulent-like fluctuations into the near-wall region faster than the mean flow can equilibrate. The result is a transient but severe amplification of aerodynamic heating, posing a significant challenge for \gls{TPS} design and with potential for localised material failure. Understanding and predicting this response remains critical for hypersonic vehicle design across a wide range of flight regimes.

Second-mode instabilities are widely recognised as the dominant mechanism in high-Mach-number boundary layers. Early controlled-forcing experiments using harmonic point sources confirmed the second mode as the primary instability and provided evidence for subharmonic resonance as a breakdown route \citep{Shiplyuk2003,Bountin2008}. More recent investigations in the Purdue quiet tunnel identified fundamental resonance as the dominant transition pathway on a Mach~6 flared cone \citep{Chynoweth2014} and established amplitude thresholds for turbulent spot formation \citep{Casper2011,Casper2012}. Complementary low-enthalpy studies, such as Mach~21 nitrogen-tunnel experiments \citep{Mironov2000}, revealed four-wave parametric resonance between fundamental and harmonic modes. However, such experimental efforts remain limited by quiet-tunnel availability, perfect-gas conditions, and incomplete access to nonlinear dynamics, highlighting the necessity for high-fidelity \gls{DNS} to systematically assess thermochemical nonequilibrium effects and wall thermal variations.

Temporal and spatial \gls{DNS} studies have investigated fundamental transition mechanisms in compressible hypersonic boundary layers. For instance, \citet{Mayer2011} computed Mach~3 transition triggered by oblique breakdown, a pathway first identified by \citet{Fasel1993}. Such simulations typically employ either broadband forcing across multiple frequencies or selective excitation of dominant energy modes \citep{AlamSandham2000,Bhaskaran2011}. Whilst the latter strategy does not mimic realistic flight or noisy wind tunnels, it provides an efficient pathway for data collection and valuable insight into underlying physics. The fundamental and subharmonic resonance mechanisms originate from the amplification of a two-dimensional instability wave. Once this wave reaches finite amplitude, secondary oblique instabilities emerge, growing either at the same frequency (fundamental resonance) or at half the frequency (subharmonic resonance). The current study aims to understand the transition mechanism through the second excitation method, building on \Cref{chapter:disturbance_growth}.

Fundamental \gls{DNS} studies of turbulent boundary-layer behaviour in hypersonic flows were initiated by \citet{Martin2007} and advanced by \citet{Duan2011}, who examined the influence of wall temperature and Mach number on mean flow statistics in cold hypersonic flows. \citet{Franko2013} performed a spatial \gls{DNS} of a Mach~6 \gls{ZPG} boundary layer, examining the skin-friction and heat-transfer overshoot observed in experiments, indicating that first-mode oblique breakdown provided a plausible transition path. An extensive database of cooled boundary layers up to Mach~14 was later generated by \citet{Zhangetal2018}.

Far fewer studies have incorporated real gas effects. \citet{Stemmer2005} simulated hypersonic boundary layer breakdown for reactive and non-reactive flow in a Mach~20 boundary layer. \citet{Marxenetal2013,Marxenetal2014} incorporated fundamental resonance and nonlinear modes in a Mach~10 adiabatic flat-plate flow, concluding that chemical reactions had no significant effect on growth rates. \citet{DiRenzo2021} performed a complete transition study incorporating chemical nonequilibrium, observing good qualitative agreement with previous works. \citet{Passiatoreetal2021} provided a comprehensive analysis of breakdown to turbulence under subharmonic second-mode excitation, observing that fundamental mode saturation and energy transfer to oblique subharmonic modes lead to $\Lambda$-vortex formation. Whilst finite-rate chemistry had little effect on initial transition stages, the endothermic nature of later stages drained modal energy, delaying transition by approximately 7\%.

Several \gls{DNS} studies have incorporated real-gas effects to examine their impact on hypersonic boundary layers, though the majority focused on chemical nonequilibrium. \citet{DiRenzo2021} investigated transition in a Mach~10 hypersonic boundary layer at suborbital enthalpy, capturing the classical sequence of second-mode amplification, resonance with oblique harmonics, and breakdown, leading to four- to five-fold increases in wall friction and heating. That work also considered finite-rate chemistry in cooled turbulent boundary layers, whilst \citet{Passiatoreetal2021} extended the framework to pseudo-adiabatic walls to isolate chemical nonequilibrium effects. A limitation of both studies was the neglect of thermal nonequilibrium. More recently, \citet{Williamsetal2025} examined chemistry--turbulence coupling, showing that turbulence drives variations in chemical source terms whilst having minimal impact on turbulence structure. Most studies have been limited to fully thermochemical equilibrium conditions, often triggered at unrealistically high Mach numbers, leaving understanding of isolated thermal effects limited.

The following chapter extends the \gls{LST} analysis in thermally excited states to the breakdown stage of hypersonic boundary layers across \gls{CPG}, \gls{TEQ}, and \gls{TNE} formulations. All cases are treated as chemically frozen, as the selected conditions do not trigger significant chemical activity \citep{BitterShepherd2015}. At high-enthalpy conditions, deviations from equilibrium become important, making it necessary to assess how different thermal models affect breakdown onset and progression. Attention is given to the interplay between linear amplification, nonlinear resonance, secondary instability growth, and turbulent spot formation. Particular emphasis is placed on the overshoot behaviour of skin friction and heat transfer, the influence of cold-wall conditions in accelerating breakdown, and the collapse of Mack second modes into fully developed turbulence. The analysis further examines the wall-temperature response of \gls{TNE} flows, where vibrational nonequilibrium modifies breakdown dynamics compared to \gls{CPG} and \gls{TEQ} cases. Collectively, these results provide a detailed characterisation of the nonlinear transition regime under different thermal models and wall conditions, highlighting the consequences for aerothermodynamic design and the challenges in developing predictive tools capable of capturing the complexity of breakdown to turbulence.


 % in Hypersonic Boundary Layers

% The development of turbulence in hypersonic boundary layers arises through a complex sequence of processes beginning with the linear growth of instability modes and culminating in nonlinear breakdown. While the second Mack mode is well captured by \gls{LST} during its early amplification, the later nonlinear dynamics that drive complete transition remain difficult to predict with comparable accuracy. The breakdown of coherent instability structures into turbulence through these mechanisms has far-reaching engineering consequences, dictating the magnitude of heating, frictional drag, and the aerothermal performance envelope of hypersonic vehicles. As disturbance amplitudes grow beyond the regime described by \gls{LST}, the boundary layer enters the final stages of transition. At this point, three-dimensionality and nonlinear effects, including secondary instabilities, dominate the dynamics and promote rapid amplification of disturbances. These interactions drive the breakdown and randomisation of initially coherent structures, ultimately establishing a fully turbulent state. Despite decades of research into nonlinear breakdown, the precise mechanisms underlying the last steps of transition remain only partially resolved. 

% The development of turbulence in hypersonic boundary layers arises through a complex sequence of processes, beginning with the linear growth of instability modes and culminating in nonlinear breakdown. Whilst the second Mack mode is well captured by \gls{LST} during its early amplification, the subsequent nonlinear dynamics that drive complete transition remain difficult to predict with comparable accuracy. The breakdown of coherent instability structures into turbulence through these mechanisms carries far-reaching engineering consequences, determining the magnitude of heating, frictional drag, and the aerothermal performance envelope of hypersonic vehicles. As disturbance amplitudes grow beyond the regime described by \gls{LST}, the boundary layer enters the final stages of transition. At this juncture, three-dimensionality and nonlinear effects, including secondary instabilities, dominate the dynamics and promote rapid amplification of disturbances. These interactions drive the breakdown and randomisation of initially coherent structures, ultimately establishing a fully turbulent state. Despite decades of research into nonlinear breakdown, the precise mechanisms underlying the ultimate steps of transition remain only partially understood. A defining feature of this regime is the overshoot in surface quantities such as the skin-friction coefficient and Stanton number. Rather than the gradual increase predicted by linear theory, both have been observed experimentally and numerically to exceed their fully turbulent levels during transition before relaxing downstream \citep{Holden1972}. This behaviour stems from the abrupt, localised enhancement of mixing and momentum transport during vortex breakdown, which injects turbulent-like fluctuations into the near-wall region faster than the mean flow can equilibrate. The result is a transient but severe amplification of aerodynamic heating, posing a significant challenge for \gls{TPS} design and with potential for localised material failure.

% The development of turbulence in hypersonic boundary layers arises through a complex sequence of processes, beginning with the linear growth of instability modes and culminating in nonlinear breakdown. Whilst the second Mack mode is well captured by \gls{LST} during its early amplification, the subsequent nonlinear dynamics that drive complete transition remain difficult to predict with comparable accuracy. The breakdown of coherent instability structures into turbulence through these mechanisms carries far-reaching engineering consequences, determining the magnitude of heating, frictional drag, and the aerothermal performance envelope of hypersonic vehicles. As disturbance amplitudes grow beyond the regime described by \gls{LST}, the boundary layer enters the final stages of transition. At this juncture, three-dimensionality and nonlinear effects, including secondary instabilities, dominate the dynamics and promote rapid amplification of disturbances. These interactions drive the breakdown and randomisation of initially coherent structures, ultimately establishing a fully turbulent state. Despite decades of research into nonlinear breakdown, the precise mechanisms underlying the ultimate steps of transition remain only partially understood. A defining feature of this regime is the overshoot in surface quantities such as the skin-friction coefficient and Stanton number. Rather than the gradual increase predicted by linear theory, both have been observed experimentally and numerically to exceed their fully turbulent levels during transition before relaxing downstream \citep{Holden1972}. This behaviour stems from the abrupt, localised enhancement of mixing and momentum transport during vortex breakdown, which injects turbulent-like fluctuations into the near-wall region faster than the mean flow can equilibrate. The result is a transient but severe amplification of aerodynamic heating, posing a significant challenge for \gls{TPS} design and with potential for localised material failure. Understanding and predicting this transient heating response remains a critical priority for the design and operation of hypersonic vehicles across a wide range of flight regimes and atmospheric conditions.


% Second-mode instabilities are widely recognised as the dominant mechanism in high-Mach-number boundary layers, and their growth has been extensively compared with linear stability theory predictions in both quiet tunnels \citep{Wilkinson1997,Chynoweth2014} and conventional noisy facilities \citep{Stetson1992}. To probe transition pathways beyond linear amplification, early controlled-forcing experiments used harmonic point sources to generate prescribed wavetrains; studies on a Mach~5.95 sharp cone confirmed the second mode as the primary instability and provided evidence for subharmonic resonance as a breakdown route \citep{Shiplyuk2003,Bountin2008}. More recent investigations in the Purdue quiet tunnel advanced this understanding under low-disturbance conditions, with \citet{Chynoweth2014} identifying fundamental resonance as the dominant transition pathway on a Mach~6 flared cone, and \citet{Casper2011,Casper2012} establishing amplitude thresholds for the onset of nonlinearity and turbulent spot formation. Complementary experiments at low enthalpy, such as the Mach~21 nitrogen-tunnel studies of \citet{Mironov2000}, showed that unsteady vortex perturbations in the shock layer can induce density fluctuations resembling natural hypersonic disturbances and revealed evidence of four-wave parametric resonance between fundamental and harmonic modes. Although these efforts have substantially advanced knowledge of transition mechanisms, they remain limited by quiet-tunnel availability, perfect-gas conditions, and incomplete access to nonlinear dynamics, highlighting the necessity for high-fidelity \gls{DNS} to systematically assess thermochemical nonequilibrium effects and wall thermal variations beyond the reach of ground-based facilities.

% A number of studies have employed temporal and spatial \gls{DNS} to investigate fundamental transition mechanisms in compressible supersonic and hypersonic boundary layers. For example, \citet{Mayer2011} carried out detailed computations of Mach~3 transition triggered by oblique breakdown, a pathway first identified by \citet{Fasel1993}. Usually, such simulations approach the transition problem in one of two ways. The first involves broadband forcing across multiple frequencies, deliberately exciting a spectrum of instability waves. The second focuses on selectively exciting dominant energy modes, as demonstrated in studies by \citet{AlamSandham2000,Bhaskaran2011}. Whilst the latter strategy does not mimic the conditions of realistic flight or noisy wind tunnels, it provides an efficient pathway for data collection and offers valuable insight into the underlying physics. The fundamental and subharmonic resonance mechanisms originate from the amplification of a two-dimensional instability wave. In hypersonic boundary layers, this primary disturbance is typically the Mack second mode, since it exhibits the highest growth rates. Once the wave reaches finite amplitude, secondary oblique instabilities emerge, growing either at the same frequency (fundamental resonance) or at half the frequency (subharmonic resonance). The current study aims to understand the transition mechanism through the second excitation method, building on the work of \Cref{chapter:disturbance_growth}.

% Fundamental \gls{DNS} studies on turbulent boundary-layer behaviour in hypersonic flows were initiated by the temporal simulations of \citet{Martin2007} and later advanced by \citet{Duan2011}, who examined the influence of wall temperature and Mach number on mean flow and turbulence statistics in cold hypersonic flows, providing a useful baseline for the present work. \citet{Franko2013} performed a spatial \gls{DNS} of a Mach~6 \gls{ZPG} boundary layer, examining the skin-friction and heat-transfer overshoot observed in experiments. Their results indicated that first-mode oblique breakdown provided a plausible path to transition, whereas fundamental second-mode excitation did not reproduce the experimental behaviour. An extensive database of cooled boundary layers up to Mach~14, with flow conditions representative of several hypervelocity wind tunnels, was later generated by \citet{Zhangetal2018}.

% Far fewer studies have included real gas effects. Stemmer~\cite{stemmer2005} conducted simulations of hypersonic boundary layer breakdown for reactive and non-reactive flow, capturing linear disturbance growth within a Mach~20 boundary layer on an isothermal flat plate without considering the nonlinear mechanisms. Marxen et al.~\cite{marxen2013,marxen2014} incorporated fundamental resonance along with nonlinear modes by exciting a finite-amplitude, two-dimensional primary wave in a Mach~10 adiabatic flat-plate flow. The study concluded that chemical reactions had no significant effect, as the chemically frozen and reacting flows exhibited identical growth rates. A \gls{DNS} study of the complete transition process, similar to that observed in \cite{franko_lele_2013}, but for chemical nonequilibrium, was performed by Di Renzo and Urzay~\cite{direnzo_urzay_2021}, who observed good qualitative agreement with previous studies. Recent studies by Passiatore et al.~\cite{passiatore2024} provided a comprehensive analysis of breakdown to turbulence for the same case as Marxen et al.~\cite{marxen2013} under subharmonic second-mode breakdown. They observed that the fundamental mode saturates quickly and transfers energy to 2D and oblique subharmonic modes, whose interactions lead to the formation of $\Lambda$-vortices. Finite-rate chemistry had little effect on the initial stages of transition. However, the endothermic nature of the later stages drained modal energy, delaying transition by approximately 7\%.

% In addition to these cold-flow investigations, several \gls{DNS} studies have incorporated real-gas effects to examine their impact on hypersonic boundary layers, although the majority of the focus was on the chemical nonequilibrium regime. Studies by \citet{Marxenetal2013,Marxenetal2014} investigated linear and nonlinear instability mechanisms in chemically active boundary layers, highlighting differences between perfect-gas, chemical nonequilibrium, and local chemical equilibrium cases, and demonstrating their influence on transition dynamics without capturing the final breakdown to turbulence. Transition in a Mach~10 hypersonic boundary layer at suborbital enthalpy was investigated by \citet{DiRenzo2021}, who captured the classical sequence of second-mode amplification, resonance with oblique harmonics, and breakdown, leading to four- to five-fold increases in wall friction and heating. That work also considered the influence of wall temperature and finite-rate chemistry in cooled turbulent boundary layers, whilst \citet{Passiatoreetal2021} extended the framework to pseudo-adiabatic walls to further isolate chemical nonequilibrium effects. A limitation of both studies was the neglect of thermal nonequilibrium, since the chosen edge conditions were not expected to induce significant vibrational excitation. More recently, the role of chemistry--turbulence coupling was examined by \citet{Williamsetal2025}, showing that turbulence drives variations in chemical source terms, whereas the impact of chemistry on the underlying turbulence structure was minimal. Most studies have been limited to fully thermochemical equilibrium conditions, often triggered at unrealistically high Mach numbers within the atmosphere. As a result, understanding of isolated thermal effects remains limited.

% The following chapter extends the analysis of \gls{LST} in thermally excited states presented previously, shifting the focus to the breakdown stage of hypersonic boundary layers across \gls{CPG}, \gls{TEQ}, and \gls{TNE} formulations. All cases are treated as chemically frozen, as the selected conditions do not trigger significant chemical activity \citep{Bitter2015}. At lower altitudes with higher densities, frequent molecular collisions maintain the gas in equilibrium \citep{ZhangZhang2022}, but at the high-enthalpy conditions relevant here, deviations from equilibrium become important, making it necessary to assess how different thermal models affect the onset and progression of breakdown. Attention is given to the interplay between linear amplification, nonlinear resonance, the growth of secondary instabilities, and the subsequent formation of turbulent spots. Particular emphasis is placed on the overshoot behaviour of skin friction and heat transfer, the influence of cold-wall conditions in accelerating breakdown, and the collapse of Mack second modes into fully developed turbulence. The analysis further examines the wall-temperature response of \gls{TNE} flows, where vibrational nonequilibrium modifies the breakdown dynamics compared to \gls{CPG} and \gls{TEQ} cases. Collectively, these results provide a detailed characterisation of the nonlinear transition regime under different thermal models and wall conditions, highlighting the consequences for aerothermodynamic design and the challenges that remain in developing predictive tools capable of capturing the complexity of breakdown to turbulence.





% A defining feature of this regime is the overshoot observed in surface quantities such as the skin-friction coefficient and Stanton number. In contrast to the gradual increase predicted by linear theory, both quantities have been shown experimentally and numerically to exceed their fully turbulent levels during transition before relaxing downstream \citep{Holden1972}. This overshoot arises from the abrupt and localised enhancement of mixing and momentum transport accompanying vortex breakdown, which deposits turbulent-like fluctuations into the near-wall region faster than the mean flow can equilibrate. The result is a transient but severe amplification of aerodynamic heating, posing a critical challenge for \gls{TPS} design and carrying the potential for localised material failure. In hypersonic applications, where baseline heat fluxes are already extreme, predicting the magnitude and streamwise location of this overshoot is therefore of paramount importance.
% A defining feature of this regime is the overshoot in surface quantities such as the skin-friction coefficient and Stanton number. Rather than the gradual increase predicted by linear theory, both have been observed experimentally and numerically to exceed their fully turbulent levels during transition before relaxing downstream \citep{Holden1972}. This behaviour stems from the abrupt, localised enhancement of mixing and momentum transport during vortex breakdown, which injects turbulent-like fluctuations into the near-wall region faster than the mean flow can equilibrate. The result is a transient but severe amplification of aerodynamic heating, posing a major challenge for \gls{TPS} design and with potential for localised material failure.

% In hypersonic applications, where baseline heat fluxes are already extreme, predicting the magnitude and location of this overshoot is therefore critical.

% Second-mode instabilities are widely recognised as the dominant mechanism in high-Mach-number boundary layers. Their growth has been examined extensively by comparing experimentally measured disturbance amplitudes with predictions from linear stability theory, both in quiet tunnels \citep{Wilkinson1997,Chynoweth2014} and in conventional noisy facilities \citep{Stetson1992}. Beyond linear amplification, a large body of experimental work has sought to clarify the mechanisms governing hypersonic transition under both flight and wind-tunnel conditions. In realistic environments, natural transition is typically initiated by broadband freestream disturbances, which in conventional hypersonic tunnels often originate from acoustic radiation of the turbulent boundary layer along the nozzle walls \citep{Schneider2015}. Early systematic investigations of hypersonic transition employed controlled forcing, where harmonic point sources were used to generate prescribed wavetrains. Experiments of this type by \citet{Shiplyuk2003} and \citet{Bountin2008} on a Mach~5.95 sharp cone confirmed the two-dimensional second mode as the primary instability and provided evidence for subharmonic resonance as a potential breakdown pathway. Building on this foundation, studies in the Purdue quiet tunnel offered further insight into natural transition under low-disturbance conditions. \citet{Chynoweth2014} examined a Mach~6 flared cone and identified fundamental resonance as the dominant route to breakdown, while \citet{Casper2011,Casper2012} investigated the growth of wave packets and the emergence of turbulent spots on the quiet-tunnel nozzle wall, establishing amplitude thresholds associated with the onset of nonlinearity and breakdown. Despite these advances, such experiments depend on specialised quiet-tunnel facilities and remain restricted in their ability to capture the full nonlinear dynamics of breakdown. Moreover, all existing hypersonic wind tunnels operate under near-perfect-gas conditions, preventing direct investigation of thermally excited flows relevant to high-enthalpy flight. \cite{Mironov2000} performed experimental studies on Mach~21 flight conditions in the hypersonic nitrogen wind tunnel T-327 of
% ITAM SB RAS, showing that unsteady vortex perturbations in the shock layer can induce density fluctuations consistent with natural hypersonic disturbances and revealed evidence of four-wave parametric resonance between fundamental and harmonic modes. These limitations underscore the need for high-fidelity \gls{DNS}, which can both complement experimental observations and systematically examine the effects of thermochemical nonequilibrium and wall thermal variations that are inaccessible in ground-based facilities.

% Second-mode instabilities are widely recognised as the dominant mechanism in high-Mach-number boundary layers, and their growth has been extensively compared with predictions from linear stability theory in both quiet tunnels \citep{Wilkinson1997,Chynoweth2014} and conventional noisy facilities \citep{Stetson1992}. To clarify transition pathways beyond linear amplification, early controlled-forcing experiments employed harmonic point sources to generate prescribed wavetrains; studies on a Mach~5.95 sharp cone confirmed the second mode as the primary instability and provided evidence for subharmonic resonance as a breakdown route \citep{Shiplyuk2003,Bountin2008}. More recent investigations in the Purdue quiet tunnel further advanced understanding under low-disturbance conditions, with \citet{Chynoweth2014} identifying fundamental resonance as the dominant transition pathway on a Mach~6 flared cone, and \citet{Casper2011,Casper2012} establishing amplitude thresholds for the onset of nonlinearity and turbulent spot formation. Complementary experiments at low enthalpy, such as the Mach~21 nitrogen-tunnel studies of \citet{Mironov2000}, demonstrated that unsteady vortex perturbations in the shock layer can induce density fluctuations similar to natural hypersonic disturbances and revealed evidence of four-wave parametric resonance between fundamental and harmonic modes. While these efforts have significantly advanced the understanding of transition mechanisms, they remain constrained by quiet-tunnel availability, perfect-gas conditions, and limited access to nonlinear dynamics, underscoring the need for high-fidelity \gls{DNS} to systematically examine thermochemical nonequilibrium effects and wall thermal variations inaccessible to ground-based facilities.

% Second-mode instabilities are widely recognised as the dominant mechanism in high-Mach-number boundary layers, and their growth has been extensively compared with linear stability theory predictions in both quiet tunnels \citep{Wilkinson1997,Chynoweth2014} and conventional noisy facilities \citep{Stetson1992}. To probe transition pathways beyond linear amplification, early controlled-forcing experiments used harmonic point sources to generate prescribed wavetrains; studies on a Mach~5.95 sharp cone confirmed the second mode as the primary instability and provided evidence for subharmonic resonance as a breakdown route \citep{Shiplyuk2003,Bountin2008}. More recent investigations in the Purdue quiet tunnel advanced this understanding under low-disturbance conditions, with \citet{Chynoweth2014} identifying fundamental resonance as the dominant transition pathway on a Mach~6 flared cone, and \citet{Casper2011,Casper2012} establishing amplitude thresholds for the onset of nonlinearity and turbulent spot formation. Complementary experiments at low enthalpy, such as the Mach~21 nitrogen-tunnel studies of \citet{Mironov2000}, showed that unsteady vortex perturbations in the shock layer can induce density fluctuations resembling natural hypersonic disturbances and revealed evidence of four-wave parametric resonance between fundamental and harmonic modes. Although these efforts have substantially advanced knowledge of transition mechanisms, they remain limited by quiet-tunnel availability, perfect-gas conditions, and incomplete access to nonlinear dynamics, highlighting the need for high-fidelity \gls{DNS} to systematically assess thermochemical nonequilibrium effects and wall thermal variations beyond the reach of ground-based facilities.


% A number of studies have employed temporal and spatial \gls{DNS} to investigate fundamental transition mechanisms in compressible supersonic and hypersonic boundary layers. For example, \citet{Mayer2011} carried out detailed computations of Mach~3 transition triggered by oblique breakdown, a pathway first identified by \citet{Fasel1993}. Usually, such simulations approach the transition problem in one of two ways. The first involves broadband forcing across multiple frequencies, deliberately exciting a spectrum of instability waves. The second focuses on selectively exciting dominant energy modes, as demonstrated in studies by \citet{AlamSandham2000,Bhaskaran2011}. While the latter strategy does not mimic the conditions of realistic flight or noisy wind tunnels, it provides an efficient pathway for data collection and offers valuable insight into the underlying physics. The fundamental and subharmonic resonance mechanisms originate from the amplification of a two-dimensional instability wave. In hypersonic boundary layers, this primary disturbance is typically the Mack second mode, since it exhibits the highest growth rates. Once the wave reaches finite amplitude, secondary oblique instabilities emerge, growing either at the same frequency (fundamental resonance) or at half the frequency (subharmonic resonance). The current studies aims to understand the transition mechanism through the second excitation method building on the works of the previous chapter. 


% Fundamental \gls{DNS} studies on turbulent boundary-layer behaviour in hypersonic flows were initiated by the temporal simulations of \citet{Martin2007} and later advanced by \citet{Duan2011}, who examined the influence of wall temperature and Mach number on mean flow and turbulence statistics in cold hypersonic flows, providing a useful baseline for the present work. \citet{Franko2013} performed a spatial \gls{DNS} of a Mach~6 \gls{ZPG} boundary layer, examining the skin-friction and heat-transfer overshoot observed in experiments. Their results indicated that first-mode oblique breakdown provided a plausible path to transition, whereas fundamental second-mode excitation did not reproduce the experimental behaviour. An extensive database of cooled boundary layers up to Mach~14, with flow conditions representative of several hypervelocity wind tunnels, was later generated by \citet{Zhangetal2018}. 

% Far fewer studies have included real gas effects. Stemmer~\cite{stemmer2005} conducted simulations of hypersonic boundary layer breakdown, for reactive and non-reactive flow, capturing linear disturbance growth within a Mach 20 boundary layer on an isothermal flat plate without considering the nonlinear mechanisms. Marxen et al. \cite{marxen2013,marxen2014} incorporated fundamental resonance along with nonlinear modes by exciting a finite-amplitude, two-dimensional primary wave in a Mach 10 adiabatic flat-plate flow. The study concluded that chemical reactions had no significant effect, as the chemically frozen and reacting flows exhibited identical growth rates. A DNS study of the complete transition process, similar to the one observed in \cite{franko_lele_2013}, but for chemical nonequilibrium, was performed by Di Renzo and Urzay~\cite{direnzo_urzay_2021}, who observed good qualitative agreement with previous studies. Recent studies by Passiatore et al.~\cite{passiatore2024} provided a comprehensive analysis of breakdown to turbulence for the same case as Marxen et. al.\cite{marxen2013} under subharmonic second-mode breakdown. They observed that the fundamental mode saturates quickly and transfers energy to 2D and oblique subharmonic modes, whose interactions lead to the formation of $\Lambda$-vortices. Finite-rate chemistry had little effect on the initial stages of transition. Still, the endothermic nature of the later stages of the flow drained modal energy, delaying transition by about 7\%.

% In addition to these cold-flow investigations, several \gls{DNS} studies have incorporated real-gas effects to examine their impact on hypersonic boundary layers, although majority of the focus was on chemical nonequilibrium portion. Studies by \citet{Marxenetal2013,Marxenetal2014} investigated linear and nonlinear instability mechanisms in chemically active boundary layers, highlighting differences between perfect-gas, chemical nonequilibrium, and local chemical equilibrium cases, and demonstrating their influence on transition dynamics without capturing the final breakdown to turbulence. Transition in a Mach~10 hypersonic boundary layer at suborbital enthalpy was investigated by \citet{DiRenzo2021}, who captured the classical sequence of second-mode amplification, resonance with oblique harmonics, and breakdown, leading to four- to five-fold increases in wall friction and heating. That work also considered the influence of wall temperature and finite-rate chemistry in cooled turbulent boundary layers, while \citet{Passiatoreetal2021} extended the framework to pseudo-adiabatic walls to further isolate chemical nonequilibrium effects. A limitation of both studies was the neglect of thermal nonequilibrium, since the chosen edge conditions were not expected to induce significant vibrational excitation. More recently, the role of chemistry–turbulence coupling was examined by \citet{Williamsetal2025}, showing that turbulence drives variations in chemical source terms, whereas the impact of chemistry on the underlying turbulence structure was minimal. .Most studies have been limited to fully thermochemical equilibrium conditions, often triggered at unrealistically high Mach numbers within the atmosphere. As a result, the understanding of isolated thermal effects remains limited. 

% The following chapter extends the analysis of \gls{LST} in thermally excited states presented previously, shifting the focus to the breakdown stage of hypersonic boundary layers across \gls{CPG}, \gls{TEQ}, and \gls{TNE} formulations. All cases are treated as chemically frozen, as the selected conditions do not trigger significant chemical activity \citep{Bitter2015}. At lower altitudes with higher densities, frequent molecular collisions maintain the gas in equilibrium \citep{ZhangZhang2022}, but at the high-enthalpy conditions relevant here, deviations from equilibrium become important, making it necessary to assess how different thermal models affect the onset and progression of breakdown. Attention is given to the interplay between linear amplification, nonlinear resonance, the growth of secondary instabilities, and the subsequent formation of turbulent spots. Particular emphasis is placed on the overshoot behaviour of skin friction and heat transfer, the influence of cold-wall conditions in accelerating breakdown, and the collapse of Mack second modes into fully developed turbulence. The analysis further examines the wall-temperature response of \gls{TNE} flows, where vibrational nonequilibrium modifies the breakdown dynamics compared to \gls{CPG} and \gls{TEQ} cases. Collectively, these results provide a detailed characterisation of the nonlinear transition regime under different thermal models and wall conditions, highlighting the consequences for aerothermodynamic design and the challenges that remain in developing predictive tools capable of capturing the complexity of breakdown to turbulence.

% \section{Numerical setup and flow definitions}
% The present study considers spatially developing, zero-pressure-gradient flat-plate boundary layers at freestream Mach numbers of $M_\infty=6$ and $M_\infty=8$. The flow conditions are based on the ICAO standard atmosphere at an altitude of $36,\text{km}$, corresponding to $p_\infty=498.5,\text{Pa}$, $T_\infty=239.3,\text{K}$, and a molar composition of $X_{N_2}=0.79$, $X_{O_2}=0.21$, yielding a stagnation enthalpy of $H=h_\infty+u_\infty^2/2=5.07,\text{MJ/kg}$. The governing equations are the compressible Navier–Stokes equations for a five-species air mixture in the chemically frozen limit.

% The computational domain is a rectangular box enclosing the boundary layer behind the shock. For the Mach 6 cases, the streamwise extent is $800\delta^*_0$, while for the Mach 8 cases it is doubled to $1600\delta^*_0$. The wall-normal and spanwise dimensions are kept fixed at $100\delta^*_0$ and $20\delta^*_0$, respectively. The corresponding grid resolutions for each case, together with the wall-to-freestream temperature ratios, are summarised in Table~\ref{tab:domain_grid}.

% \begin{table}[hbt!]
% \centering
% \caption{Grid discretisation for the four simulated cases.}
% \begin{tabular}{lcccc}
% \toprule
% Case & $M_\infty$ & $T_w/T_\infty$ & Grid resolution ($N_x \times N_y \times N_z$) & Total points \\
% \midrule
% M6Tw12 & 6 & 12 & $5000 \times 600 \times 305$ & $9.15\times 10^8$ \\
% M6Tw18 & 6 & 18 & $5200 \times 600 \times 275$ & $8.58\times 10^8$ \\
% M8Tw12 & 8 & 12 & $4600 \times 500 \times 160$ & $3.68\times 10^8$ \\
% M8Tw18 & 8 & 18 & $4600 \times 500 \times 160$ & $3.68\times 10^8$ \\
% \bottomrule
% \end{tabular}
% \label{tab:domain_grid}
% \end{table}



% The grid is uniform in the streamwise, $x$ and spanwise $z$ directions, with the wall-normal spacing generated through a hyperbolic tangent stretching to cluster points near the wall as $y = L_y \sinh(s_1 \xi)/\sinh(s_1)$ for uniformly distributed points on $\xi \in [0,1]$, with a stretching factor $s_1 = 5.0$. Inflow boundary-layer profiles are obtained from a locally self-similar framework solved using Chebyshev collocation, allowing consistent specification across \gls{CPG}, \gls{TEQ}, and \gls{TNE} conditions, with the \gls{TNE} case employing an induced \gls{TFG} similarity solution to provide thermally frozen inflow states. Disturbances are introduced downstream of the inflow through a wall-normal suction–blowing forcing function with Gaussian streamwise localisation and multi-frequency content designed to excite multiple fundamental instabilities and examine their nonlinear interaction. The wall is non-catalytic and isothermal, with two distinct temperature conditions investigated: a cold wall at $T_w=1200 \ \text{K}$ and a hot wall at $T_w=1800 \ \text{K}$. These wall conditions are essential for assessing the sensitivity of the boundary layer to realistic surface cooling, where cold walls enhance the amplification of high-frequency second-mode instabilities by thinning the boundary layer and shifting the stability spectrum, accelerating breakdown and amplifying overshoot in skin friction and heat transfer, while hot walls act to stabilise acoustic-trapped disturbances and delay transition. By systematically varying both Mach number and wall temperature, the simulations provide a framework to examine their combined influence on instability amplification, nonlinear breakdown, and the development of turbulence. As mentioned previously, the conservative variables are initialised by solving the similarity solution which is transferred into the domain through an inlet pressure extrapolation boundary condition. The outflow and far-field boundary conditions are treated using simple extrapolation and zero-gradient boundary conditions respectively. Periodic boundary conditions are imposed in the spanwise direction, and the inflow Reynolds number is set to $Re_{\delta^*_0}=10{,}000$, ensuring comparability across cases and sufficient domain length to capture the full transition sequence.




% \section{Numerical setup and flow definitions}

% The present study considers spatially developing, zero-pressure-gradient flat-plate boundary layers at freestream Mach numbers of $M_\infty=6$ and $M_\infty=8$. Flow conditions are based on the ICAO standard atmosphere at an altitude of 36~km, corresponding to $p_\infty=498.5~\text{Pa}$, $T_\infty=239.3~\text{K}$, and a molar composition of $X_{N_2}=0.79$ and $X_{O_2}=0.21$, yielding a stagnation enthalpy of $H=h_\infty+u_\infty^2/2=5.07~\text{MJ/kg}$. The governing equations are the compressible Navier--Stokes equations for a five-species air mixture in the chemically frozen limit.

% The computational domain is a rectangular box enclosing the boundary layer, with streamwise extent of $800\delta^*_0$ for Mach~6 cases and $1600\delta^*_0$ for Mach~8 cases. The wall-normal and spanwise dimensions are fixed at $100\delta^*_0$ and $20\delta^*_0$, respectively. The grid is uniform in the streamwise and spanwise directions, with wall-normal spacing generated through hyperbolic tangent stretching to cluster points near the wall as $y = L_y \sinh(s_1 \xi)/\sinh(s_1)$ for uniformly distributed points on $\xi \in [0,1]$, with a stretching factor $s_1 = 5.0$. Grid resolutions for each case are summarised in Table~\ref{tab:domain_grid}.

% \begin{table}[hbt!]
% \centering
% \caption{Grid discretisation and wall temperature ratios for the four simulated cases.}
% \begin{tabular}{lcccc}
% \toprule
% Case & $M_\infty$ & $T_w/T_\infty$ & Grid resolution ($N_x \times N_y \times N_z$) & Total points \\
% \midrule
% M6Tw12 & 6 & 12 & $5000 \times 600 \times 305$ & $9.15\times 10^8$ \\
% M6Tw18 & 6 & 18 & $5200 \times 600 \times 275$ & $8.58\times 10^8$ \\
% M8Tw12 & 8 & 12 & $4600 \times 500 \times 160$ & $3.68\times 10^8$ \\
% M8Tw18 & 8 & 18 & $4600 \times 500 \times 160$ & $3.68\times 10^8$ \\
% \bottomrule
% \end{tabular}
% \label{tab:domain_grid}
% \end{table}

% Inflow boundary-layer profiles are obtained from a locally self-similar framework solved using Chebyshev collocation, enabling consistent specification across \gls{CPG}, \gls{TEQ}, and \gls{TNE} conditions, with the \gls{TNE} case employing an induced \gls{TFG} similarity solution to provide thermally frozen inflow states. Conservative variables are initialised by solving the similarity solution and transferred into the domain through an inlet pressure extrapolation boundary condition. The outflow and far-field boundary conditions employ simple extrapolation condition, whilst periodic boundary conditions are imposed in the spanwise direction. Disturbances are introduced downstream of the inflow through a wall-normal suction--blowing forcing function with Gaussian streamwise localisation and multi-frequency content designed to excite multiple fundamental instabilities and examine their nonlinear interaction. The wall is non-catalytic and isothermal, with two distinct temperature conditions investigated: a cold wall at $T_w=1200~\text{K}$ and a hot wall at $T_w=1800~\text{K}$. Cold walls enhance amplification of high-frequency second-mode instabilities by thinning the boundary layer and shifting the stability spectrum, thereby accelerating breakdown and amplifying overshoot in skin friction and heat transfer. Hot walls act to stabilise acoustic-trapped disturbances and delay transition. By systematically varying both Mach number and wall temperature, the simulations provide a framework to examine their combined influence on instability amplification, nonlinear breakdown, and the development of turbulence. The inflow Reynolds number is set to $Re_{\delta^*_0}=10{,}000$, ensuring comparability across cases and sufficient domain length to capture the full transition sequence.

% \subsection{Disturbance seeding and excitation method}

% To induce transition to turbulence, a suction--blowing forcing strip based on the approach of \cite{AlamSandham2000} is employed, extended here to a two-dimensional domain and modified for large-amplitude perturbations. A Gaussian-shaped disturbance is imposed at $x_f$, as illustrated in~\cref{figure:schematicflow}, and introduced through a wall mass-flux perturbation designed to promote nonlinear breakdown. The forcing is prescribed as
% \begin{equation}
% v_w(x,z,t) \;=\; a_0 \, g(x) \, h(z) \, \sum_{i=1}^{n} \sin(2\pi f_i t),
% \end{equation}
% where the streamwise envelope is
% \begin{equation}
% g(x) = \exp\left(-\frac{(x-x_f)^2}{2(\delta^*_0)^2}\right),
% \end{equation}
% and the spanwise modulation is
% \begin{equation}
% h(z) = 1 + \sin\!\left(\frac{2\pi z}{105\delta^*_0} + \frac{2\pi z}{205\delta^*_0}\right).
% \end{equation} where the Gaussian envelope of width defined by $2\sigma^2 = 20(\delta^*_0)^2$, centred at $x_f = 25\delta^*_0$ from the origin of the simulation domain. This placement ensures sufficient grid resolution within the shape function to accurately represent the disturbance. The forcing amplitude is set to $a_0 = 0.1 \ u_\infty$, chosen to be sufficiently large to trigger transition within the computational domain while remaining representative of realistic disturbance levels. The frequency set is case dependent. For Mach~6, $\{f_i\} = \{0.15,\,0.16\}$ is applied, while for Mach~8, $\{f_i\} = \{0.10,\,0.11\}$. These frequencies are selected based on local linear stability theory to ensure that the second-mode growth reaches its maximum between $x/\delta^*_0 \approx 220$ and $280$, coinciding with the forcing location at $x_f \approx 240\delta^*_0$ as mentioned in the previous chapter. This alignment ensures that the forcing excites the naturally most amplified disturbances and produces sustained growth within the intended region of the computational domain.

% \begin{figure}[t]
%     \centering
%     \includegraphics[width=10cm]{resources/figures/breakdown/insert_shapefunction_natural.pdf}
%     \caption{Shape function for the Breakdown forcing.}
%     \label{figure:forcing_shapefunction}
% \end{figure}

% The frequency set is case dependent. For Mach~6, $\{f_i\} = \{0.15,\,0.16\}$ is applied, while for Mach~8, $\{f_i\} = \{0.10,\,0.11\}$. These frequencies are selected based on local linear stability theory to ensure that the second-mode growth reaches its maximum between $x/\delta^*_0 \approx 220$ and $280$, coinciding with the forcing location at $x_f \approx 240\delta^*_0$. This alignment ensures that the forcing excites the naturally most amplified disturbances and produces sustained growth within the intended region of the computational domain. Promoting nonlinear interactions between harmonics and subharmonics is known to accelerate breakdown. No oblique modes are imposed, in contrast to the classical oblique-transition route ~\citep{Franko2013,Williamsetal2025}, in order to isolate the evolution of two-dimensional second-mode disturbances and their intrinsic secondary instabilities. Instead, three-dimensionalisation is introduced through the spanwise modulation term, which imposes long- and short-wavelength subharmonics ($\lambda_z=105\delta^*_0$ and $205\delta^*_0$), allowing $\Lambda$-shaped vortices and streaks to emerge naturally without biasing the flow towards an artificially prescribed oblique structure.

% \begin{figure}[t]
%     \centering
%     \includegraphics[width=12cm]{resources/figures/breakdown/result_mach6_stability_diagram.pdf}
%     \caption{Stability plot at $\delta_{*}^{0} = 255$ for (a) CPG, (b) TEQ, (c) TNE.}
%     \label{figure:stability_plot}
% \end{figure}

% The simulations are advanced in time with non-dimensional timesteps of $\Delta t = 0.012$ for the Mach~6 cases and $\Delta t = 0.016$ for the Mach~8 cases. Each computation is performed for two flow-through times, defined as $t_f = L_x/u_\infty$, after which the solution is restarted and continued for a further two flow-throughs during which Favre-averaged statistics are accumulated. Owing to the high-frequency nature of the induced acoustic waves, the transition dynamics evolve rapidly and the total integration time required to capture breakdown is relatively modest. This procedure removes the initial transients associated with inflow forcing and wavepacket adjustment, ensuring that the collected statistics represent the fully developed transitional and turbulent state. The resulting dataset thus captures the complete transition sequence, from the controlled excitation of second-mode instabilities through nonlinear harmonic and subharmonic interactions to the onset of turbulence.

\section{Numerical setup and flow definitions}

This study examines spatially developing, zero-pressure-gradient flat-plate boundary layers at freestream Mach numbers of $M_\infty=6$ and $M_\infty=8$. Flow conditions correspond to the ICAO standard atmosphere at 36~km altitude: $p_\infty=498.5~\text{Pa}$, $T_\infty=239.3~\text{K}$, and a five-species air composition ($X_{N_2}=0.79$, $X_{O_2}=0.21$) yielding a stagnation enthalpy of $H=h_\infty+u_\infty^2/2=5.07~\text{MJ/kg}$. The governing equations are the compressible Navier--Stokes equations in the chemically frozen limit. The computational domain is a rectangular box with streamwise extents of $800\delta^*_0$ (Mach~6) and $1600\delta^*_0$ (Mach~8), wall-normal height of $100\delta^*_0$, and spanwise width of $20\delta^*_0$. The grid is uniform in streamwise and spanwise directions, with wall-normal spacing generated via hyperbolic tangent stretching as $y = L_y \sinh(s_1 \xi)/\sinh(s_1)$ for uniformly distributed $\xi \in [0,1]$ and stretching factor $s_1 = 5.0$. Grid resolutions are summarised in Tables~\ref{tab:domain_grid_dns,tab:domain_grid_iles}.

% \begin{table}[hbt!]
% \centering
% \caption{Grid resolution and thermal conditions for all simulated cases.}
% \begin{tabular}{lcccc}
% \toprule
% Case & $M_\infty$ & $T_w/T_\infty$ & Grid resolution ($N_x \times N_y \times N_z$) & Total points \\
% \midrule
% M6Tw12 & 6 & 12 & $5000 \times 600 \times 305$ & $9.15\times 10^8$ \\
% M6Tw18 & 6 & 18 & $5200 \times 600 \times 275$ & $8.58\times 10^8$ \\
% M8Tw12 & 8 & 12 & $4600 \times 500 \times 160$ & $3.68\times 10^8$ \\
% M8Tw18 & 8 & 18 & $4600 \times 500 \times 160$ & $3.68\times 10^8$ \\
% \bottomrule
% \end{tabular}
% \label{tab:domain_grid}
% \end{table}


\begin{table}[hbt!]
\centering
\caption{Grid resolution and thermal conditions for DNS simulations.}
\begin{tabular}{lccccc}
\toprule
Case & Gas Model & $M_\infty$ & $T_w/T_\infty$ &  ($N_x \times N_y \times N_z$) & Total points \\
\midrule
\multirow{3}{*}{M6Tw18} & CPG & 6 & 18 & $5200 \times 600 \times 275$ & $8.58\times 10^8$ \\
 & TEQ & 6 & 18 & $5200 \times 600 \times 275$ & $8.58\times 10^8$ \\
 & TNE & 6 & 18 & $5200 \times 600 \times 275$ & $8.58\times 10^8$ \\
\bottomrule
\end{tabular}
\label{tab:domain_grid_dns}
\end{table}

\begin{table}[hbt!]
\centering
\caption{Grid resolution and thermal conditions for ILES simulations.}
\begin{tabular}{lccccc}
\toprule
Case & Gas Model & $M_\infty$ & $T_w/T_\infty$ &  ($N_x \times N_y \times N_z$) & Total points \\
\midrule
M6Tw12 & TNE & 6 & 12 & $2600 \times 600 \times 205$ & $3.20\times 10^8$ \\
M6Tw18 & TNE & 6 & 18 & $2300 \times 500 \times 165$ & $1.84\times 10^8$ \\
M8Tw12 & TNE & 8 & 12 & $4600 \times 500 \times 165$ & $3.68\times 10^8$ \\
M8Tw18 & TNE & 8 & 18 & $4600 \times 500 \times 165$ & $3.68\times 10^8$ \\
\bottomrule
\end{tabular}
\label{tab:domain_grid_iles}
\end{table}

Inflow boundary-layer profiles are obtained from a locally self-similar framework solved using Chebyshev collocation, enabling consistent specification across \gls{CPG}, \gls{TEQ}, and \gls{TNE} conditions. For the \gls{TNE} case, an induced \gls{TFG} similarity solution provides thermally frozen inflow states. Conservative variables are initialised by solving the similarity equations and imposed via inlet pressure extrapolation. The outflow and far-field employ simple extrapolation conditions, whilst spanwise boundaries are periodic.

The wall is non-catalytic and isothermal at two temperatures: $T_w=1200~\text{K}$ (cold wall) or $T_w=1800~\text{K}$ (hot wall). Cold walls enhance high-frequency second-mode amplification by thinning the boundary layer and shifting the stability spectrum, thereby accelerating breakdown and amplifying skin-friction and heat-transfer overshoot. Hot walls stabilise acoustic-trapped disturbances and delay transition. This systematic variation examines the combined effects of Mach number and wall temperature on instability amplification, nonlinear breakdown, and turbulence development.

The inflow Reynolds number is set to $Re_{\delta^*_0}=10{,}000$, ensuring comparability across cases and sufficient domain length to capture the full transition sequence. The governing equations are discretised using the WENO-Z scheme for convective terms to suppress non-physical oscillations, combined with 4th order central differences for viscous and metric terms. Time integration is performed using a 3rd order strong stability preserving Runge--Kutta scheme. Transport properties are evaluated using the Musawi--Sandham viscosity model and vibrational relaxation is calculated using the Millikan--White formulation with the Park correction. Simulations are advanced with non-dimensional timesteps $\Delta t = 0.012$ (Mach~6) and $\Delta t = 0.016$ (Mach~8), which maintain stability whilst resolving the high-frequency acoustic waves characteristic of the second-mode instability.



\subsection{Disturbance seeding and excitation method}

Transition to turbulence is induced via a suction--blowing forcing strip, following the approach of \cite{AlamSandham2000} and extended here to a two-dimensional domain with large-amplitude perturbations. A Gaussian-shaped disturbance is imposed at $x_f$ through a wall mass-flux perturbation designed to promote nonlinear breakdown:
\begin{equation}
v_w(x,z,t) \;=\; a_0 \, g(x) \, h(z) \, \sum_{i=1}^{n} \sin(2\pi f_i t),
\end{equation}
where the streamwise envelope is
\begin{equation}
g(x) = \exp\left(-\frac{(x-x_f)^2}{2(\delta^*_0)^2}\right),
\end{equation}
and the spanwise modulation is
\begin{equation}
h(z) = 1 + \sin\!\left(\frac{2\pi z}{105\delta^*_0} + \frac{2\pi z}{205\delta^*_0}\right).
\end{equation}
The Gaussian envelope has width $2\sigma^2 = 20(\delta^*_0)^2$ and is centred at $x_f = 25\delta^*_0$ from the domain origin, ensuring sufficient grid resolution to accurately represent the disturbance. The forcing amplitude is set to $a_0 = 0.1 \, u_\infty$, sufficiently large to trigger transition within the domain whilst remaining representative of realistic disturbance levels. The frequency set is case-dependent: for Mach~6, $\{f_i\} = \{0.15,\,0.16\}$, and for Mach~8, $\{f_i\} = \{0.10,\,0.11\}$. These frequencies are selected using local linear stability theory to ensure that second-mode growth reaches its maximum between $x/\delta^*_0 \approx 220$ and $280$, coinciding with the forcing location at $x_f \approx 240\delta^*_0$. This alignment ensures the forcing excites the most naturally amplified disturbances and sustains growth within the intended computational region.

\begin{figure}[t]
    \centering
    \includegraphics[width=10cm]{resources/figures/breakdown/insert_shapefunction_natural.pdf}
    \caption{Spatial structure of the wall mass-flux forcing function.}
    \label{figure:forcing_shapefunction}
\end{figure}

Rather than imposing oblique modes (as in classical oblique-transition routes), three-dimensionality is introduced through spanwise modulation, which imposes long- and short-wavelength subharmonics ($\lambda_z=105\delta^*_0$ and $205\delta^*_0$). This approach is justified by the stability analysis shown in Figure~\ref{figure:stability_plot}, which demonstrates the absence of first-mode instability across all thermal models; consequently, oblique modes which arise from first-mode interactions need not be artificially introduced. Instead, $\Lambda$-shaped vortices and streaks emerge naturally, isolating the evolution of two-dimensional second-mode disturbances and their secondary instabilities without prescribing an artificial oblique structure.

\begin{figure}[t]
    \centering
    \includegraphics[width=12cm]{resources/figures/breakdown/result_mach6_stability_diagram.pdf}
    \caption{Linear stability analysis at $\delta_{*}^{0} = 255$ for (a) CPG, (b) TEQ, and (c) TNE models.}
    \label{figure:stability_plot}
\end{figure}

Each simulation spans four consecutive flow-through times ($t_f = L_x/u_\infty$). The first two flow-throughs remove initial transients from inflow forcing and wavepacket adjustment; the solution is then restarted and continued for two further flow-throughs whilst accumulating Favre-averaged statistics. This procedure ensures that statistics represent the fully developed transitional state, isolated from transient effects. The total integration time is compact due to the high-frequency nature of induced acoustic waves and rapid transition dynamics. The resulting dataset captures the complete sequence from controlled second-mode excitation through nonlinear harmonic and subharmonic interactions to turbulence onset.

% \section{Numerical setup and flow definitions}

% This study examines spatially developing, zero-pressure-gradient flat-plate boundary layers at freestream Mach numbers of $M_\infty=6$ and $M_\infty=8$. Flow conditions correspond to the ICAO standard atmosphere at 36~km altitude: $p_\infty=498.5~\text{Pa}$, $T_\infty=239.3~\text{K}$, and a five-species air composition ($X_{N_2}=0.79$, $X_{O_2}=0.21$) yielding a stagnation enthalpy of $H=h_\infty+u_\infty^2/2=5.07~\text{MJ/kg}$. The governing equations are the compressible Navier--Stokes equations in the chemically frozen limit. The computational domain is a rectangular box with streamwise extents of $800\delta^*_0$ (Mach~6) and $1600\delta^*_0$ (Mach~8), wall-normal height of $100\delta^*_0$, and spanwise width of $20\delta^*_0$. The grid is uniform in streamwise and spanwise directions, with wall-normal spacing generated via hyperbolic tangent stretching as $y = L_y \sinh(s_1 \xi)/\sinh(s_1)$ for uniformly distributed $\xi \in [0,1]$ and stretching factor $s_1 = 5.0$. Grid resolutions are summarised in Table~\ref{tab:domain_grid}.

% \begin{table}[hbt!]
% \centering
% \caption{Grid resolution and thermal conditions for all simulated cases.}
% \begin{tabular}{lcccc}
% \toprule
% Case & $M_\infty$ & $T_w/T_\infty$ & Grid resolution ($N_x \times N_y \times N_z$) & Total points \\
% \midrule
% M6Tw12 & 6 & 12 & $5000 \times 600 \times 305$ & $9.15\times 10^8$ \\
% M6Tw18 & 6 & 18 & $5200 \times 600 \times 275$ & $8.58\times 10^8$ \\
% M8Tw12 & 8 & 12 & $4600 \times 500 \times 160$ & $3.68\times 10^8$ \\
% M8Tw18 & 8 & 18 & $4600 \times 500 \times 160$ & $3.68\times 10^8$ \\
% \bottomrule
% \end{tabular}
% \label{tab:domain_grid}
% \end{table}

% Inflow boundary-layer profiles are obtained from a locally self-similar framework solved using Chebyshev collocation, enabling consistent specification across \gls{CPG}, \gls{TEQ}, and \gls{TNE} conditions. For the \gls{TNE} case, an induced \gls{TFG} similarity solution provides thermally frozen inflow states. Conservative variables are initialised by solving the similarity equations and imposed via inlet pressure extrapolation. The outflow and far-field employ simple extrapolation conditions, whilst spanwise boundaries are periodic.

% The wall is non-catalytic and isothermal at two temperatures: $T_w=1200~\text{K}$ (cold wall) or $T_w=1800~\text{K}$ (hot wall). Cold walls enhance high-frequency second-mode amplification by thinning the boundary layer and shifting the stability spectrum, thereby accelerating breakdown and amplifying skin-friction and heat-transfer overshoot. Hot walls stabilise acoustic-trapped disturbances and delay transition. This systematic variation examines the combined effects of Mach number and wall temperature on instability amplification, nonlinear breakdown, and turbulence development.

% Disturbances are introduced via a wall-normal suction--blowing forcing function with Gaussian streamwise localisation and multi-frequency content, designed to excite multiple fundamental instabilities and examine their nonlinear interactions. The inflow Reynolds number is set to $Re_{\delta^*_0}=10{,}000$, ensuring comparability across cases and sufficient domain length to capture the full transition sequence.


% \subsection{Disturbance seeding and excitation method}

% Transition to turbulence is induced via a suction--blowing forcing strip, following the approach of \cite{AlamSandham2000} and extended here to a two-dimensional domain with large-amplitude perturbations. A Gaussian-shaped disturbance is imposed at $x_f$ through a wall mass-flux perturbation designed to promote nonlinear breakdown:
% \begin{equation}
% v_w(x,z,t) \;=\; a_0 \, g(x) \, h(z) \, \sum_{i=1}^{n} \sin(2\pi f_i t),
% \end{equation}
% where the streamwise envelope is
% \begin{equation}
% g(x) = \exp\left(-\frac{(x-x_f)^2}{2(\delta^*_0)^2}\right),
% \end{equation}
% and the spanwise modulation is
% \begin{equation}
% h(z) = 1 + \sin\!\left(\frac{2\pi z}{105\delta^*_0} + \frac{2\pi z}{205\delta^*_0}\right).
% \end{equation}
% The Gaussian envelope has width $2\sigma^2 = 20(\delta^*_0)^2$ and is centred at $x_f = 25\delta^*_0$ from the domain origin, ensuring sufficient grid resolution to accurately represent the disturbance. The forcing amplitude is set to $a_0 = 0.1 \, u_\infty$, sufficiently large to trigger transition within the domain whilst remaining representative of realistic disturbance levels.

% The frequency set is case-dependent: for Mach~6, $\{f_i\} = \{0.15,\,0.16\}$, and for Mach~8, $\{f_i\} = \{0.10,\,0.11\}$. These frequencies are selected using local linear stability theory to ensure that second-mode growth reaches its maximum between $x/\delta^*_0 \approx 220$ and $280$, coinciding with the forcing location at $x_f \approx 240\delta^*_0$. This alignment ensures the forcing excites the most naturally amplified disturbances and sustains growth within the intended computational region.

% \begin{figure}[t]
%     \centering
%     \includegraphics[width=10cm]{resources/figures/breakdown/insert_shapefunction_natural.pdf}
%     \caption{Spatial structure of the wall mass-flux forcing function.}
%     \label{figure:forcing_shapefunction}
% \end{figure}

% Rather than imposing oblique modes (as in classical oblique-transition routes), three-dimensionalisation is introduced through spanwise modulation, which imposes long- and short-wavelength subharmonics ($\lambda_z=105\delta^*_0$ and $205\delta^*_0$). This approach allows $\Lambda$-shaped vortices and streaks to emerge naturally, isolating the evolution of two-dimensional second-mode disturbances and their secondary instabilities without prescribing an artificial oblique structure.

% \begin{figure}[t]
%     \centering
%     \includegraphics[width=12cm]{resources/figures/breakdown/result_mach6_stability_diagram.pdf}
%     \caption{Linear stability analysis at $\delta_{*}^{0} = 255$ for (a) CPG, (b) TEQ, and (c) TNE models.}
%     \label{figure:stability_plot}
% \end{figure}

% Simulations are advanced with non-dimensional timesteps $\Delta t = 0.012$ (Mach~6) and $\Delta t = 0.016$ (Mach~8). Each computation spans two flow-through times ($t_f = L_x/u_\infty$), after which the solution is restarted and continued for two further flow-throughs whilst accumulating Favre-averaged statistics. This procedure removes initial transients from inflow forcing and wavepacket adjustment, ensuring statistics represent the fully developed transitional state. The total integration time is modest due to the high-frequency nature of induced acoustic waves and rapid transition dynamics. The resulting dataset captures the complete sequence from controlled second-mode excitation through nonlinear harmonic and subharmonic interactions to turbulence onset.



% \citep{DiRenzo2021} studied the transition of Mach~10 flow hypersonic boundary layer at suborbital enthalpy to examine transition, turbulence, and air dissociation. Transition followed a classical path via two-dimensional second-mode amplification, resonance with oblique harmonics, and eventual breakdown, leading to four- to five-fold increases in wall friction and heating. The results show that transition mechanisms remain robust at high enthalpy, but highlight open challenges involving vibrational non-equilibrium, catalytic and rough-wall effects, wall response, pressure gradients, and geometry.
% The interaction of finite-rate chemistry with transitional and turbulent spatially developing boundary layers was recently investigated by \citet{DiRenzo2021}, who analysed the effects of chemical nonequilibrium in cooled turbulent boundary layers. Subsequent work by \citet{Passiatoreetal2021} focused on pseudo-adiabatic wall conditions, isolating the role of chemical nonequilibrium under those constraints. In both cases, thermal nonequilibrium was neglected, as the selected edge conditions were not expected to induce significant vibrational relaxation.

% In the Purdue quiet tunnel, \citet{Chynoweth2014} investigated natural transition on a Mach~6 flared cone and identified fundamental resonance as the dominant route to breakdown. Complementary experiments by \citet{Casper2011,Casper2012} analysed the development of wave packets and turbulent spots on the quiet-tunnel nozzle wall, establishing amplitude thresholds associated with the onset of nonlinearity and breakdown. To probe the transition process in a more systematic and repeatable manner, controlled forcing has also been employed, in which harmonic point sources generate prescribed wavetrains. Experiments of this type by ~\cite{Shiplyuk2003} and ~\cite{Bountin2008} on a Mach~5.95 sharp cone confirmed the two-dimensional second mode as the leading instability and provided evidence for subharmonic resonance as a viable breakdown pathway.


% 
% Second-mode instabilities are generally regarded as the dominant instability mechanism in high-Mach-number boundary layers. Their growth has been studied in detail through comparisons of experimentally measured disturbance amplitudes with predictions from linear stability theory, both in quiet tunnels \citep{Wilkinson1997,Chynoweth2014} and in conventional noisy facilities \citep{Stetson1992}. A considerable body of experimental work has been devoted to clarifying the mechanisms of hypersonic transition under both flight and wind-tunnel conditions. In realistic environments, natural transition is typically triggered by broadband freestream disturbances, which in conventional hypersonic tunnels often originate from acoustic radiation of the turbulent boundary layer along the nozzle walls \citep{Schneider2015}. ~\cite{Chynoweth2014} experimentally studied natural transition on a Mach~6 flared cone in the Purdue quiet tunnel and identified fundamental resonance as the dominant path to turbulence breakdown. ~\cite{Casper2011,Casper2012} investigated the development of wave packets and turbulent spots in a hypersonic boundary layer on the quiet-tunnel nozzle wall, and reported amplitude thresholds associated with the onset of nonlinearity and breakdown.

%  Complementary studies by Casper et al.~\citep{Casper2011,Casper2012} tracked the growth of wave packets and the formation of turbulent spots on the quiet-tunnel nozzle wall, establishing amplitude thresholds associated with the onset of nonlinearity and breakdown. To probe the transition process in a more systematic and repeatable manner, controlled forcing has also been employed, in which harmonic point sources generate prescribed wavetrains. Experiments of this type by Shiplyuk et al.~\citep{16} and Bountin et al.~\citep{17} on a Mach~5.95 sharp cone confirmed the two-dimensional second mode as the leading instability and provided evidence for subharmonic resonance as a viable breakdown pathway.



% The progression from a laminar to a turbulent boundary layer in hypersonic flows represents one of the most critical yet least predictable phenomena in high-speed aerothermodynamics. Although the early amplification of instability waves---most prominently the Mack second mode---can be accurately characterised within the framework of linear stability theory, the later stages of transition are dominated by nonlinear mechanisms that fall outside the scope of linear analysis. 

% Understanding and modelling this regime is therefore of central importance for the design of thermal protection systems, where even small shifts in transition onset or the intensity of breakdown can translate into significant changes in local heat flux.  rea











% \section{Influence of thermal excitations on flow behaviour}
% \subsection{Transitional regime}

% \section{Influence of thermal excitations on flow behaviour}

% The influence of thermal modelling is examined by comparing \gls{CPG}, \gls{TEQ}, and \gls{TNE} formulations for the Mach~6 hot-wall case. Previous linear stability results demonstrate that \gls{TEQ} exhibits the highest second-mode growth rates, exceeding those of both \gls{TNE} and \gls{CPG}. This arises from stronger coupling of vibrationally equilibrated modes with acoustic disturbances under \gls{TEQ}, whilst \gls{TNE} introduces vibrational relaxation that suppresses amplification and \gls{CPG} yields the weakest growth due to the absence of thermal excitation. These distinctions establish the context for interpreting the nonlinear evolution in the simulations.

% % \begin{figure}[t]
% %     \centering
% %     \includegraphics[width=6.3cm]{resources/figures/test/rhou1_contour.pdf}
% %     \caption{Linear stability analysis at $\delta_{*}^{0} = 255$ for (a) CPG, (b) TEQ, and (c) TNE models.}
% %     \label{figure:stability_plot}
% % \end{figure}


% Transition is characterised using statistical diagnostics and flow-field decomposition in frequency and spanwise wavenumber space to identify dominant instability mechanisms and their role in producing heat-transfer overshoot. Transition onset is determined from the sharp rise in skin friction and Stanton number. Mean profiles, velocity fluctuations, turbulence intensities, spanwise correlations, and energy spectra are examined to trace the progression from laminar to turbulent flow and reveal the development of three-dimensional structures. Reynolds stresses and wall heat fluxes are evaluated within transitional and turbulent regions to clarify energy redistribution and the mechanisms sustaining turbulence, with particular attention to overshoot phenomena.


\section{Influence of thermal excitations on flow behaviour}

The influence of thermal modelling assumptions on hypersonic boundary layer behaviour is investigated through systematic comparison of \gls{CPG}, \gls{TEQ}, and \gls{TNE} formulations for the Mach~6 hot-wall configuration. Linear stability analyses established in preceding work reveal that thermal modelling exerts a profound influence on the second-mode growth characteristics. Specifically, \gls{TEQ} exhibits the highest second-mode growth rates, substantially exceeding those predicted by both \gls{TNE} and \gls{CPG} approaches. This enhanced amplification arises from the stronger modal coupling between vibrationally equilibrated energy modes and acoustic disturbances within the \gls{TEQ} framework. In contrast, the \gls{TNE} formulation introduces vibrational relaxation effects that attenuate the amplification of these modes, thereby reducing the overall growth rates. The \gls{CPG} model, which neglects thermal excitation entirely, yields the weakest instability growth, attributable to the absence of the additional energy channels that characterise excited-state molecular motion. These distinctions in linear stability behaviour provide essential context for interpreting the subsequent nonlinear evolution observed within the direct numerical simulations and for understanding how thermal nonequilibrium effects manifest in the transitional flow dynamics.

The characterisation of transition is accomplished through comprehensive analysis of statistical diagnostics and flow-field quantities throughout the computational domain. The onset of transition is determined objectively from the sharp rise in both the skin-friction coefficient and Stanton number, which serve as robust indicators of the laminar-to-turbulent transformation. The temporal and spatial evolution is traced through careful examination of mean velocity profiles, velocity fluctuation components, turbulence intensity distributions, and spanwise two-point correlations, which collectively illuminate the progression from laminar flow through the transitional regime to fully developed turbulence, whilst revealing the development of coherent three-dimensional structures characteristic of the nonlinear breakdown process and their role in heat-transfer enhancement.



% \newpage 

\begin{figure}[t]
    \centering
    \includegraphics[width=18.9cm,angle=-90]{resources/figures/test/result_contour_dns_skinfriction.pdf}
    \caption{Linear stability analysis at $\delta_{*}^{0} = 255$ for (a) CPG, (b) TEQ, and (c) TNE models.}
    \label{figure:stability_plot}
\end{figure}

The energy redistribution mechanisms sustaining the transitional and turbulent flow states are clarified through detailed evaluation of Reynolds stress components and wall-normal heat flux distributions within both the transitional and fully turbulent regions of the domain. Special emphasis is placed on quantifying the mechanisms responsible for the characteristic heat-transfer overshoot relative to the asymptotic turbulent state, identifying whether such phenomena arise from localised coherent-structure interactions, transient amplification mechanisms, or residual effects of the weakly nonlinear breakdown process. These analyses, conducted systematically across the three thermal models, illuminate the extent to which vibrational non-equilibrium and thermal excitation effects modify the fundamental turbulence-production mechanisms and alter the scaling of wall-bounded transport phenomena in hypersonic flows.

% \newpage

% \subsection{Early disturbance amplification}

% The streamwise evolution of wall quantities characterises the onset and progression of transition, as shown in Figure~\ref{figure:thermal_cfst}. The skin-friction coefficient and Stanton number are the primary diagnostics:
% \begin{equation}
%     C_f = \cfrac{2\tau_w}{\rho_\infty u_\infty^2},
% \end{equation}
% \begin{equation}
%     St = \cfrac{q_w}{\rho_\infty u_\infty c_{p,tr} (T_{aw}-T_w)},
% \end{equation}
% where $T_{aw}$ is the adiabatic wall temperature obtained from the self-similar solution of the corresponding flow conditions. Since the recovery equation does not yield accurate temperatures for \gls{TEQ} and \gls{TNE} cases, this approach ensures consistency across all thermal models. The Stanton number is evaluated using the translational--rotational temperature $T$ for \gls{CPG} and \gls{TNE}, whilst for \gls{TEQ} it employs the combined $T$ and $T_v$. The skin-friction coefficient reflects wall shear stress and the growth of near-wall disturbances and turbulent streaks. The Stanton number quantifies nondimensional wall heat flux, indicating the thermal response of the boundary layer and heat-transfer efficiency. Together, these parameters reveal momentum--energy coupling during transition, with their streamwise evolution typically exhibiting overshoot relative to the turbulent asymptote, a signature of nonlinear breakdown. For laminar boundary layers, the skin-friction and Stanton number coefficients are obtained directly from similarity solutions. Following \cite{White2005}, these are expressed as  
% \begin{equation}
% c_{f,lam} = \sqrt{2} f''(0) \,Re_x^{-1/2}\cfrac{\mu_w \rho_w}{\mu_\infty \rho_\infty},
% \end{equation}
% \begin{equation}
% St_{lam} = \cfrac{\theta'(0)Re_x^{-1/2}}{\sqrt{2Pr_{tr}\,(T_r-T_w)}}\cfrac{\mu_w \rho_w}{\mu_\infty \rho_\infty},
% \end{equation}
% where $f$ and $\theta$ are the similarity solutions for velocity and temperature, and $T_r$ is the recovery factor obtained from the similarity solution. Note that both expressions vary across thermal models; the formulations shown employ the TNE model, with analogous expressions for \gls{CPG} and \gls{TEQ} arising from their respective similarity solutions.

% For turbulent boundary layers, predictions are compared against the compressible Van Driest~II correlation \citep{Sandhametal2014}:
% \begin{equation}
%     c_{f,turb} = \frac{1}{\lambda} \left[ \frac{0.242 \left( \arcsin\alpha + \arcsin\beta \right)}{\log_{10}(c_f Re_x) - 0.76 \log_{10}(T_w/T_\infty) + 0.41} \right],
% \end{equation}
% where $\lambda = (\gamma - 1) M_\infty^2 r/2$, $\alpha = (2a^2 - b)/Q$, $\beta = b/Q$, $Q = \sqrt{b^2 + 4a^2}$, $a = \sqrt{\lambda T_\infty/T_w}$, and $b = (1+\lambda)T_\infty/T_w - 1$. This implicit equation converges rapidly with a suitable initial guess. The corresponding turbulent Stanton number is obtained via the Reynolds analogy:
% \begin{equation}
% St_{turb} = \frac{c_f}{2Pr_{tr}^{2/3}}.
% \end{equation}

\subsection{Early disturbance amplification}

The streamwise evolution of wall-bounded quantities provides essential diagnostics for characterising the onset and progression of boundary layer transition, as illustrated in Figure~\ref{figure:thermal_cfst}. Two primary parameters are employed: the skin-friction coefficient and Stanton number, defined respectively as
\begin{equation}
    C_f = \frac{2\tau_w}{\rho_\infty u_\infty^2},
\end{equation}
\begin{equation}
    St = \frac{q_w}{\rho_\infty u_\infty c_{p,tr} (T_{aw}-T_w)},
\end{equation}
where $T_{aw}$ denotes the adiabatic wall temperature derived from the self-similar solution corresponding to the specified flow conditions. The adiabatic wall temperature is calculated directly from the similarity solution rather than via the recovery equation, as the latter yields insufficient accuracy for \gls{TEQ} and \gls{TNE} regimes. This approach ensures methodological consistency across all thermal models. The Stanton number is evaluated using the translational--rotational temperature $T$ for both \gls{CPG} and \gls{TNE} cases, whereas the \gls{TEQ} formulation employs a weighted combination of $T$ and $T_v$. The skin-friction coefficient reflects the wall shear stress and the development of near-wall disturbances and turbulent streaks, whilst the Stanton number quantifies the nondimensional heat flux at the wall, indicating the thermal response of the boundary layer and heat-transfer efficiency. The streamwise evolution of these parameters reveals the coupled momentum--energy dynamics during transition, characteristically exhibiting overshoot relative to the fully turbulent asymptote, which is a hallmark signature of nonlinear breakdown.

For laminar boundary layers, the skin-friction and Stanton number coefficients are derived directly from self similarity solutions. Following \cite{White2005}, these are expressed as
\begin{equation}
c_{f,lam} = \sqrt{2} f''(0) \,Re_x^{-1/2}\frac{\mu_w \rho_w}{\mu_\infty \rho_\infty},
\end{equation}
\begin{equation}
St_{lam} = \frac{\theta'(0)Re_x^{-1/2}}{\sqrt{2Pr_{tr}\,(T_r-T_w)}}\frac{\mu_w \rho_w}{\mu_\infty \rho_\infty},
\end{equation}
where $f$ and $\theta$ represent the dimensionless velocity and temperature similarity functions respectively, and $T_r$ is the recovery temperature obtained from the corresponding similarity solution. Both expressions exhibit thermal-model dependence; the formulations presented employ the \gls{TNE} model, with analogous expressions for \gls{CPG} and \gls{TEQ} arising naturally from their respective similarity solutions.

For turbulent boundary layers, predictions are benchmarked against the compressible Van Driest~II correlation \citep{Sandhametal2014}:
\begin{equation}
    c_{f,turb} = \frac{1}{\lambda} \left[ \frac{0.242 \left( \arcsin\alpha + \arcsin\beta \right)}{\log_{10}(c_f Re_x) - 0.76 \log_{10}(T_w/T_\infty) + 0.41} \right],
\end{equation}
where the auxiliary parameters are defined as $\lambda = (\gamma - 1) M_\infty^2 r/2$, $\alpha = (2a^2 - b)/Q$, $\beta = b/Q$, $Q = \sqrt{b^2 + 4a^2}$, $a = \sqrt{\lambda T_\infty/T_w}$, and $b = (1+\lambda)T_\infty/T_w - 1$. This implicit equation is solved iteratively, converging rapidly when initialised with an appropriate first estimate. The corresponding turbulent Stanton number is subsequently obtained via the Reynolds analogy:
\begin{equation}
St_{turb} = \frac{c_f}{2Pr_{tr}^{2/3}}.
\end{equation}

\begin{figure}[t]
    \centering
    \fontsize{9}{10}\selectfont
    \includegraphics[width=0.85\textwidth]{resources/figures/test/result_wallquantities_cf_st.pdf}
    \caption[Streamwise evolution of $C_f$ and $St$ for Mach~6, $T_w=1800$ K.]%
    {Streamwise evolution of spatially temporally averaged  $C_f$ and $St$ for the Mach~6, $T_w=1800$ K case. Results are shown for \gls{CPG}, \gls{TEQ}, and \gls{TNE} formulations, alongside laminar reference solutions and the Van Driest~II turbulent correlation.}
    \label{figure:thermal_cfst}
\end{figure}



% For laminar boundary layers, the skin-friction and Stanton number coefficients are obtained directly from similarity solutions. Following \cite{White2005}, these are expressed as  
% \begin{equation}
% c_{f,lam} = \sqrt{2f''(0)}\,Re_x^{-1/2}\frac{\mu_w \rho_w}{\mu_\infty \rho_\infty},
% \qquad
% St_{lam} = \frac{g'(0)Re_x^{-1/2}}{\sqrt{2Pr_{tr}\,(T_r-T_w)}}\frac{\mu_w \rho_w}{\mu_\infty \rho_\infty},
% \end{equation}
% where $f$ and $g$ denote the similarity solutions for velocity and temperature, respectively, and $T_r$ is the recovery factor calculated through the similarity solution.  

% For turbulent boundary layers, a suitable extension to the current study is to compare against the compressible Van Driest~II correlation \citep{Sandhametal2014}, given by  
% \begin{equation}
%     c_{f,turb} = \frac{1}{\lambda} \left[ \frac{0.242 \left( \arcsin\alpha + \arcsin\beta \right)}{\log_{10}(c_f Re_x) - 0.76 \log_{10}(T_w/T_\infty) + 0.41} \right],
% \end{equation}
% with $\lambda = (\gamma - 1) M_\infty^2 r/2$, $\alpha = (2a^2 - b)/Q$, $\beta = b/Q$, $Q = \sqrt{b^2 + 4a^2}$, $a = \sqrt{\lambda T_\infty/T_w}$, and $b = (1+\lambda)T_\infty/T_w - 1$. The formulation requires an implicit iteration for $c_f$, but converges rapidly with a suitable initial guess. The corresponding Stanton number is obtained through the standard Reynolds analogy,  
% \begin{equation}
% St_{turb} = \frac{c_f}{2Pr_{tr}^{2/3}}.
% \end{equation}

% \begin{figure}[t]
%     \centering
%     \fontsize{9}{10}\selectfont
%     \includegraphics[width=0.98\textwidth]{resources/figures/breakdown/result_thermal_cf_st.pdf}
%     \caption[Streamwise evolution of $C_f$ and $St$ for Mach~6, $T_w=1800$ K.]%
%     {Streamwise evolution of spatially temporally averaged  $C_f$ and $St$ for the Mach~6, $T_w=1800$ K case. Results are shown for \gls{CPG}, \gls{TEQ}, and \gls{TNE} formulations, alongside laminar reference solutions and the Van Driest~II turbulent correlation.}
%     \label{figure:thermal_cfst}
% \end{figure}


In the present work, laminar and turbulent reference profiles are shown only for the \gls{CPG} case, to avoid excessive clutter in the plots. Comparisons indicate that the laminar and turbulent trends for \gls{CPG} and \gls{TNE} are very similar to those for \gls{TEQ}, although the \gls{TEQ} profiles lie slightly lower. Therefore, the use of \gls{CPG} profiles as representative references is sufficient.


% \begin{figure}[t]
%     \centering
%     \fontsize{9}{10}\selectfont
%     \includegraphics[width=1.04\textwidth]{resources/figures/breakdown/result_thermal_contours_st.pdf}
%     \caption[Time-averaged wall Stanton number distribution across the span for Mach~6, $T_w=1800$ K.]%
%     {Time-averaged spanwise distribution of wall Stanton number $St$ for the Mach~6, $T_w=1800$ K case. The contours reveal the development of three-dimensional thermal patterns during transition, with enhanced streak organisation in the transitional region and a more uniform distribution as the flow approaches turbulence. Results are compared across \gls{CPG}, \gls{TEQ}, and \gls{TNE} formulations.}
%     \label{figure:thermal_spanwise_st}
% \end{figure}

% \begin{figure}[t]
%     \centering
%     \fontsize{9}{10}\selectfont
%     \includegraphics[width=21.3cm,angle=-90]{resources/figures/breakdown/result_thermal_contours_st.pdf}
%     \caption[Time-averaged wall Stanton number distribution across the span for Mach~6, $T_w=1800$ K.]%
%     {Time-averaged spanwise distribution of wall Stanton number $St$ for the Mach~6, $T_w=1800$ K case, shown for \gls{CPG}, \gls{TEQ}, and \gls{TNE} formulations.}
%     \label{figure:thermal_spanwise_st}
% \end{figure}


% As shown in Figure~\ref{figure:thermal_cfst}, transition onset varies significantly across thermal models: \gls{TEQ} undergoes transition earliest, followed by \gls{CPG}, whilst \gls{TNE} exhibits the latest onset. An initial forcing spike is imposed to investigate \gls{LST} predictions, as \gls{DNS} results are known to underpredict \gls{LST} growth rates \citep{Stemmer2005}. In all cases, a small elevation appears in skin-friction and heat-transfer distributions upstream of breakdown, associated with second-mode amplification predicted by \gls{LST}. The initial laminar regions exhibit slight variations across models due to differences in the boundary-layer thermal state. This elevation reflects primary disturbance resonance of the introduced \textit{Mack modes}, after which breakdown proceeds through secondary instabilities \citep{DiRenzo2021}, consistent with observations in \gls{CPG} flows \citep{Franko2013} but absent in other studies \citep{Passiatoreetal2021,Williamsetal2025}. An overshoot relative to the turbulent asymptote follows, characteristic of nonlinear breakdown. The overshoot is most pronounced in \gls{TEQ} and smoother in \gls{TNE}, which exhibits a more gradual relaxation to the turbulent state. This sequence—\gls{TEQ} transitioning first, followed by \gls{CPG} and \gls{TNE}—highlights the stabilising effect of thermal nonequilibrium. Notably, this ordering contrasts with the $N$-factor growth of the previous chapter: although \gls{TEQ} showed the strongest amplification consistent with its early transition, the weaker growth of \gls{CPG} would suggest later transition than observed. The results confirm that \gls{TNE} provides the most effective delay of transition and moderates overshoot behaviour.

% % Selected flow quantities and grid parameters at representative stations throughout the transition process are presented in Table~\ref{table:global_quantities_models}.

% Selected flow quantities and grid parameters at representative stations throughout the transition process are presented in Table~\ref{table:global_quantities_models}. The table documents the streamwise evolution across five characteristic phases: near resonance of the two-dimensional second-mode, secondary instability growth, imminent breakdown, end of transition, and the fully turbulent state. Boundary-layer integral parameters including the momentum thickness Reynolds number $Re_\theta$, shape factor $H$, and displacement thickness Reynolds number $Re^*_\delta$ characterise the development of the velocity profile. Thermal quantities including the maximum temperature rise $\tilde{T}_{max}/T_e$ and vibrational energy budget terms $B_q^{TR}$ and $B_q^{V}$ quantify thermochemical effects, with the latter revealing the opposing signs characteristic of translational-rotational and vibrational energy balance in \gls{TNE} flows. The friction Reynolds number $Re_\tau$ and associated grid spacings $\Delta x^+$, $\Delta y_w^+$, and $\Delta z^+$ document the computational resolution requirements, which increase substantially during transition as coherent structures develop and the boundary layer thickens.

As shown in Figure~\ref{figure:thermal_cfst}, transition onset varies significantly across thermal models: \gls{TEQ} undergoes transition earliest, followed by \gls{CPG}, whilst \gls{TNE} exhibits the latest onset. An initial forcing spike is imposed to investigate \gls{LST} predictions, as \gls{DNS} results are known to underpredict \gls{LST} growth rates \citep{Stemmer2005}. In all cases, a small elevation appears in skin-friction and heat-transfer distributions upstream of breakdown, associated with second-mode amplification predicted by \gls{LST}. The initial laminar regions exhibit slight variations across models due to differences in the boundary-layer thermal state. This elevation reflects primary disturbance resonance of the introduced \textit{Mack modes}, after which breakdown proceeds through secondary instabilities \citep{DiRenzo2021}, consistent with observations in \gls{CPG} flows \citep{Franko2013} but absent in other studies \citep{Passiatoreetal2021,Williamsetal2025}. An overshoot relative to the turbulent asymptote follows, characteristic of nonlinear breakdown. The overshoot is most pronounced in \gls{TEQ} and smoother in \gls{TNE}, which exhibits a more gradual relaxation to the turbulent state. This sequence—\gls{TEQ} transitioning first, followed by \gls{CPG} and \gls{TNE}—highlights the stabilising effect of thermal nonequilibrium. Notably, this ordering contrasts with the $N$-factor growth of the previous chapter: although \gls{TEQ} showed the strongest amplification consistent with its early transition, the weaker growth of \gls{CPG} would suggest later transition than observed. The results confirm that \gls{TNE} provides the most effective delay of transition and moderates overshoot behaviour.


% \begin{figure}[t]
%     \centering
%     \fontsize{9}{10}\selectfont
%     \includegraphics[width=0.95\textwidth]{resources/figures/test/result_contour_dns_skinfriction_new.pdf}
%     \caption[Streamwise evolution of $C_f$ and $St$ for Mach~6, $T_w=1800$ K.]%
%     {Streamwise evolution of time-averaged  $C_f$ for the Mach~6, $T_w=1800$ K case. Results are shown for \gls{CPG}, \gls{TEQ}, and \gls{TNE}.}
%     \label{figure:contour_skinfriction}
% \end{figure}

Selected flow quantities and grid parameters at representative stations throughout the transition process are presented in Table~\ref{table:global_quantities_models}, documenting streamwise evolution across five characteristic phases: near resonance of the two-dimensional second-mode, secondary instability growth, imminent breakdown, end of transition, and the fully turbulent state. Boundary-layer integral parameters $Re_\theta$, $H$, and $Re^*_\delta$ characterise velocity profile development, with the shape factor decreasing from approximately $13$--$14$ in laminar and early transition regions to below $11$ in the turbulent state, reflecting progressive profile steepening as coherent structures develop. Thermal quantities including maximum temperature rise $\tilde{T}_{max}/T_e$ and vibrational energy terms $B_q^{TR}$ and $B_q^{V}$ quantify thermochemical effects; in \gls{TNE} flows, the latter reveals opposing signs characteristic of translational-rotational and vibrational energy balance. The friction Reynolds number $Re_\tau$ ranges from approximately $150$ in the resonance phase to over $800$ in the turbulent state for \gls{TNE}, reflecting increased near-wall velocity gradients during breakdown. Grid spacings $\Delta x^+$, $\Delta y_w^+$, and $\Delta z^+$ satisfy established \gls{DNS} resolution requirements, ensuring adequate capture of near-wall viscous structures and streamwise coherent features.

% Selected flow quantities and grid parameters at representative stations throughout the transition process are presented in Table~\ref{table:global_quantities_models}; the table documents the streamwise evolution across five characteristic phases: near resonance of the two-dimensional second-mode, secondary instability growth, imminent breakdown, end of transition, and the fully turbulent state. Boundary-layer integral parameters including the momentum thickness Reynolds number $Re_\theta$, shape factor $H$, and displacement thickness Reynolds number $Re^*_\delta$ characterise the development of the velocity profile throughout transition. The shape factor evolution is particularly revealing, decreasing from approximately $13$--$14$ in the laminar and early transition regions to values below $11$ in the turbulent state, reflecting the progressive steepening of the velocity profile as coherent structures develop. Thermal quantities including the maximum temperature rise $\tilde{T}_{max}/T_e$ and vibrational energy budget terms $B_q^{TR}$ and $B_q^{V}$ quantify thermochemical effects, with the latter revealing the opposing signs characteristic of translational-rotational and vibrational energy balance in \gls{TNE} flows. The friction Reynolds number $Re_\tau$ exhibits substantial variation across the transition process, ranging from approximately $150$ in the resonance phase to over $800$ in the turbulent state for \gls{TNE}, reflecting the dramatic increase in near-wall velocity gradients during breakdown. The associated grid spacings $\Delta x^+$, $\Delta y_w^+$, and $\Delta z^+$ satisfy established \gls{DNS} resolution requirements throughout, ensuring adequate capture of near-wall viscous structures and streamwise coherent features across all transition phases.






% As shown in Figure~\ref{figure:thermal_cfst}, the \gls{TEQ} case undergoes transition earliest, marked by a sharp rise ahead of the other formulations. The \gls{CPG} case follows with a slightly delayed onset, while the \gls{TNE} solutions display the latest transition. In all cases, a small bulge appears in the skin-friction and heat-transfer distributions upstream of breakdown, representing the amplification of the Mack second mode predicted by \gls{LST}. This feature is associated with an initial resonance of the primary disturbances, after which the breakdown occurs further downstream due to the action of secondary instabilities. An overshoot relative to the turbulent asymptote is then evident, characteristic of nonlinear breakdown. The overshoot is most pronounced in \gls{TEQ}, whereas \gls{TNE} exhibits a smoother approach to the turbulent state. The ordering, with \gls{TEQ} first, then \gls{CPG}, and finally \gls{TNE}, highlights the influence of thermal nonequilibrium in delaying transition and moderating overshoot behaviour. This trend contrasts with the $N$-factor growth presented in the previous chapter: while \gls{TEQ} exhibited the largest amplification, consistent with its early transition here, the \gls{CPG} case would be expected to transition later than observed owing to its weaker $N$-factor growth. Instead, the results confirm that \gls{TNE} most effectively delays the transition process. The spanwise distributions of Stanton number, shown in Figure~\ref{figure:thermal_spanwise_st}, provide further evidence of the three-dimensionalisation process during transition. In the early stages, the wall heat flux remains nearly uniform across the span, consistent with the dominance of two-dimensional second-mode disturbances. As transition progresses, streak-like structures emerge, reflecting the growth of spanwise modulations that redistribute heat flux into alternating high- and low-intensity bands. These patterns mark the onset of secondary instabilities that destabilise the initially two-dimensional Mack modes and promote rapid breakdown to turbulence. By the end of the transitional region, the streak organisation weakens and the wall heat flux becomes more evenly distributed across the span, consistent with a fully developed turbulent state.






% The streamwise evolution of wall quantities is first examined. Figures~\ref{fig:Cf}--\ref{fig:Cq} present the distributions of the skin-friction coefficient $C_f$ and the normalized wall heat fluxes $C_q$ and $C_{q,V}$, defined as
% \begin{equation}
% C_f = \frac{2\tau_w}{\rho_\infty u_\infty^2},
% \qquad
% C_q = \frac{q_w}{\rho_\infty u_\infty^3},
% \qquad
% C_{q,V} = \frac{q_{w,V}}{\rho_\infty u_\infty^3},
% \label{eq:coeffs}
% % \end{equation}
% where $\tau_w$ is the wall shear stress, $q_w$ is the total wall heat flux, and $q_{w,V}$ denotes the vibrational contribution. These coefficients provide a direct measure of the wall response during transition: $C_f$ reflects the momentum transfer through shear, while $C_q$ and $C_{q,V}$ quantify the convective and vibrational thermal loads. Their streamwise evolution highlights the overshoot in skin friction and heat transfer that accompanies nonlinear breakdown and the onset of turbulence.

\newpage

% \begin{sidewaystable}[H]
% \centering
% \caption{Streamwise evolution of selected dimensionless flow quantities for .}
% \begin{tabular}{lccc ccc ccc ccc ccc}
% \toprule
%  & \multicolumn{3}{c}{Near resonance of 2D mode} 
%  & \multicolumn{3}{c}{Secondary instability} 
%  & \multicolumn{3}{c}{Imminent breakdown} 
%  & \multicolumn{3}{c}{End of transition} 
%  & \multicolumn{3}{c}{Turbulent} \\
% \cmidrule(lr){2-4} \cmidrule(lr){5-7} \cmidrule(lr){8-10} \cmidrule(lr){11-13} \cmidrule(lr){14-16}
%  & CPG & TEQ & TNE & CPG & TEQ & TNE & CPG & TEQ & TNE & CPG & TEQ & TNE & CPG & TEQ & TNE \\
% \midrule
% $\hat{x}$ & ... & ... & ... & ... & ... & ... & ... & ... & ... & ... & ... & ... & ... & ... & ... \\
% $Re_x \times 10^{-6}$ & ... & ... & ... & ... & ... & ... & ... & ... & ... & ... & ... & ... & ... & ... & ... \\
% $Re_\tau$ & ... & ... & ... & ... & ... & ... & ... & ... & ... & ... & ... & ... & ... & ... & ... \\
% $Re^*_\delta \times 10^{-4}$ & ... & ... & ... & ... & ... & ... & ... & ... & ... & ... & ... & ... & ... & ... & ... \\
% $Re_\theta$ & ... & ... & ... & ... & ... & ... & ... & ... & ... & ... & ... & ... & ... & ... & ... \\
% $Re_\theta^{inc}$ & ... & ... & ... & ... & ... & ... & ... & ... & ... & ... & ... & ... & ... & ... & ... \\
% $H$ & ... & ... & ... & ... & ... & ... & ... & ... & ... & ... & ... & ... & ... & ... & ... \\
% $M_{a,t,max}$ & ... & ... & ... & ... & ... & ... & ... & ... & ... & ... & ... & ... & ... & ... & ... \\
% $B_q^{TR}$ & ... & ... & ... & ... & ... & ... & ... & ... & ... & ... & ... & ... & ... & ... & ... \\
% $B_q^{V}$ & ... & ... & ... & ... & ... & ... & ... & ... & ... & ... & ... & ... & ... & ... & ... \\
% $\tilde{T}_{max}/T_e$ & ... & ... & ... & ... & ... & ... & ... & ... & ... & ... & ... & ... & ... & ... & ... \\
% $\delta_{M_t}/\delta^{*}_{0}$ & ... & ... & ... & ... & ... & ... & ... & ... & ... & ... & ... & ... & ... & ... & ... \\
% $\delta_{T}/\delta^{*}_{0}$ & ... & ... & ... & ... & ... & ... & ... & ... & ... & ... & ... & ... & ... & ... & ... \\
% \midrule
% $x^+$ & ... & ... & ... & ... & ... & ... & ... & ... & ... & ... & ... & ... & ... & ... & ... \\
% $y^+$ & ... & ... & ... & ... & ... & ... & ... & ... & ... & ... & ... & ... & ... & ... & ... \\
% $z^+$ & ... & ... & ... & ... & ... & ... & ... & ... & ... & ... & ... & ... & ... & ... & ... \\
% $L_z^+$ & ... & ... & ... & ... & ... & ... & ... & ... & ... & ... & ... & ... & ... & ... & ... \\
% \bottomrule
% \end{tabular}
% \label{table:global_quantities_models}
% \end{sidewaystable}

% \newpage
% \begin{sidewaystable}[H]
% \centering
% % \small
% \caption{Streamwise evolution of selected dimensionless flow quantities for TNE simulations across the transition process.}
% \begin{tabular}{lccc ccc ccc ccc ccc}
% \toprule
%  & \multicolumn{3}{c}{Near resonance of 2D mode} 
%  & \multicolumn{3}{c}{Secondary instability} 
%  & \multicolumn{3}{c}{Imminent breakdown} 
%  & \multicolumn{3}{c}{End of transition} 
%  & \multicolumn{3}{c}{Turbulent} \\
% \cmidrule(lr){2-4} \cmidrule(lr){5-7} \cmidrule(lr){8-10} \cmidrule(lr){11-13} \cmidrule(lr){14-16}
%  & CPG & TEQ & TNE & CPG & TEQ & TNE & CPG & TEQ & TNE & CPG & TEQ & TNE & CPG & TEQ & TNE \\
% \midrule
% $x/\delta^*_0$ & ... & 252 & 292 & ... & 291 & 361 & ... & 337 & 436 & ... & 388 & 505 & ... & 600 & 600 \\
% $Re_x \times 10^{-6}$ & ... & 4.29 & 4.53 & ... & 4.68 & 5.22 & ... & 5.14 & 5.97 & ... & 5.65 & 6.66 & ... & 7.77 & 7.61 \\
% $Re_\tau$ & ... & 145 & 189 & ... & 151 & 166 & ... & 175 & 240 & ... & 254 & 514 & ... & 687 & 824 \\
% $Re^*_\delta \times 10^{-4}$ & ... & 1.56 & 1.68 & ... & 1.63 & 1.80 & ... & 1.71 & 1.93 & ... & 1.79 & 2.04 & ... & 2.10 & 2.18 \\
% $Re_\theta$ & ... & 1132 & 1211 & ... & 1200 & 1333 & ... & 1294 & 1537 & ... & 1437 & 1809 & ... & 2123 & 2323 \\
% $H$ & ... & 13.22 & 13.96 & ... & 13.11 & 14.00 & ... & 12.59 & 12.69 & ... & 11.73 & 11.03 & ... & 10.04 & 10.71 \\
% $M_{a,t,max}$ & ... & 0.08 & 0.09 & ... & 0.08 & 0.07 & ... & 0.08 & 0.09 & ... & 0.09 & 0.14 & ... & 0.11 & 0.15 \\
% $B_q^{TR} \times 10^{-2}$ & ... & $-2.93$ & $-5.11$ & ... & $-2.88$ & $-3.78$ & ... & $-3.00$ & $-4.72$ & ... & $-3.41$ & $-6.99$ & ... & $-4.21$ & $-7.23$ \\
% $B_q^{V} \times 10^{-3}$ & ... & $-5.87$ & 5.58 & ... & $-5.77$ & 5.79 & ... & $-6.01$ & 5.64 & ... & $-6.83$ & 5.69 & ... & $-8.45$ & 5.43 \\
% $\tilde{T}_{max}/T_e$ & ... & 3.36 & 3.73 & ... & 3.37 & 3.67 & ... & 3.40 & 3.75 & ... & 3.48 & 3.81 & ... & 3.51 & 3.72 \\
% $\delta_{M_t}/\delta^{*}_{0}$ & ... & 0.90 & 0.90 & ... & 0.90 & 0.90 & ... & 0.89 & 0.86 & ... & 0.87 & 0.67 & ... & 0.66 & 0.33 \\
% \midrule
% $\Delta x^+$ & ... & 5.49 & 6.56 & ... & 5.47 & 5.31 & ... & 5.88 & 6.72 & ... & 6.67 & 10.05 & ... & 8.35 & 10.97 \\
% $\Delta y_w^+$ & ... & 0.48 & 0.58 & ... & 0.48 & 0.47 & ... & 0.52 & 0.59 & ... & 0.59 & 0.88 & ... & 0.73 & 0.96 \\
% $\Delta y^+_{99}$ & ... & 1.52 & 1.97 & ... & 1.58 & 1.72 & ... & 1.81 & 2.46 & ... & 2.59 & 5.20 & ... & 6.89 & 8.27 \\
% $\Delta z^+$ & ... & 2.55 & 3.05 & ... & 2.54 & 2.47 & ... & 2.73 & 3.12 & ... & 3.10 & 4.67 & ... & 3.88 & 5.09 \\
% $L_z^+$ & ... & 714 & 853 & ... & 711 & 690 & ... & 764 & 873 & ... & 867 & 1306 & ... & 1085 & 1426 \\
% \bottomrule
% \end{tabular}
% \label{table:global_quantities_models}
% \end{sidewaystable}
% \newpage


% \newpage
% \begin{sidewaystable}[H]
% \centering
% % \small
% \caption{Streamwise evolution of selected dimensionless flow quantities for TNE simulations across the transition process.}
% \begin{tabular}{lccc ccc ccc ccc ccc}
% \toprule
%  & \multicolumn{3}{c}{Near resonance of 2D mode} 
%  & \multicolumn{3}{c}{Secondary instability} 
%  & \multicolumn{3}{c}{Imminent breakdown} 
%  & \multicolumn{3}{c}{End of transition} 
%  & \multicolumn{3}{c}{Turbulent} \\
% \cmidrule(lr){2-4} \cmidrule(lr){5-7} \cmidrule(lr){8-10} \cmidrule(lr){11-13} \cmidrule(lr){14-16}
%  & CPG & TEQ & TNE & CPG & TEQ & TNE & CPG & TEQ & TNE & CPG & TEQ & TNE & CPG & TEQ & TNE \\
% \midrule
% $x/\delta^*_0$ & 252 & 252 & 292 & 291 & 291 & 361 & 337 & 337 & 436 & 388 & 388 & 505 & 789 & 600 & 600 \\
% $Re_x \times 10^{-6}$ & 4.12 & 4.29 & 4.53 & 4.51 & 4.68 & 5.22 & 4.97 & 5.14 & 5.97 & 5.48 & 5.65 & 6.66 & 9.49 & 7.77 & 7.61 \\
% $Re_\tau$ & 214 & 145 & 189 & 193 & 151 & 166 & 359 & 175 & 240 & 1041 & 254 & 514 & 5611 & 687 & 824 \\
% $Re^*_\delta \times 10^{-4}$ & 1.60 & 1.56 & 1.68 & 1.68 & 1.63 & 1.80 & 1.76 & 1.71 & 1.93 & 1.85 & 1.79 & 2.04 & 2.43 & 2.10 & 2.18 \\
% $Re_\theta$ & 1168 & 1132 & 1211 & 1271 & 1200 & 1333 & 1498 & 1294 & 1537 & 1881 & 1437 & 1809 & 6355 & 2123 & 2323 \\
% $H$ & 13.63 & 13.22 & 13.96 & 13.48 & 13.11 & 14.00 & 11.48 & 12.59 & 12.69 & 8.78 & 11.73 & 11.03 & 4.70 & 10.04 & 10.71 \\
% $M_{a,t,max}$ & 0.10 & 0.08 & 0.09 & 0.09 & 0.08 & 0.07 & 0.14 & 0.08 & 0.09 & 0.16 & 0.09 & 0.14 & 0.14 & 0.11 & 0.15 \\
% $B_q^{TR} \times 10^{-2}$ & $-7.56$ & $-2.93$ & $-5.11$ & $-5.15$ & $-2.88$ & $-3.78$ & $-7.68$ & $-3.00$ & $-4.72$ & $-8.96$ & $-3.41$ & $-6.99$ & $-7.37$ & $-4.21$ & $-7.23$ \\
% $B_q^{V} \times 10^{-3}$ & 0.00 & $-5.87$ & 5.58 & 0.00 & $-5.77$ & 5.79 & 0.00 & $-6.01$ & 5.64 & 0.00 & $-6.83$ & 5.69 & 0.00 & $-8.45$ & 5.43 \\
% $\tilde{T}_{max}/T_e$ & 3.70 & 3.36 & 3.73 & 3.69 & 3.37 & 3.67 & 3.77 & 3.40 & 3.75 & 3.79 & 3.48 & 3.81 & 3.71 & 3.51 & 3.72 \\
% $\delta_{M_t}/\delta^{*}_{0}$ & 0.88 & 0.90 & 0.90 & 0.90 & 0.90 & 0.90 & 0.83 & 0.89 & 0.86 & 0.61 & 0.87 & 0.67 & 0.82 & 0.66 & 0.33 \\
% \midrule
% $\Delta x^+$ & 7.76 & 5.49 & 6.56 & 6.63 & 5.47 & 5.31 & 9.84 & 5.88 & 6.72 & 11.91 & 6.67 & 10.05 & 10.66 & 8.35 & 10.97 \\
% $\Delta y_w^+$ & 0.68 & 0.48 & 0.58 & 0.58 & 0.48 & 0.47 & 0.86 & 0.52 & 0.59 & 1.05 & 0.59 & 0.88 & 0.94 & 0.73 & 0.96 \\
% $\Delta y^+_{99}$ & 2.24 & 1.52 & 1.97 & 2.01 & 1.58 & 1.72 & 3.68 & 1.81 & 2.46 & 10.43 & 2.59 & 5.20 & 55.95 & 6.89 & 8.27 \\
% $\Delta z^+$ & 3.60 & 2.55 & 3.05 & 3.08 & 2.54 & 2.47 & 4.57 & 2.73 & 3.12 & 5.53 & 3.10 & 4.67 & 4.95 & 3.88 & 5.09 \\
% $L_z^+$ & 1008 & 714 & 853 & 862 & 711 & 690 & 1279 & 764 & 873 & 1548 & 867 & 1306 & 1385 & 1085 & 1426 \\
% \bottomrule
% \end{tabular}
% \label{table:global_quantities_models}
% \end{sidewaystable}
% \newpage

\newpage
\begin{sidewaystable}[H]
\centering
% \small
\caption{Streamwise evolution of selected dimensionless flow quantities for TNE simulations across the transition process.}
\begin{tabular}{lccc ccc ccc ccc ccc}
\toprule
 & \multicolumn{3}{c}{Near resonance of 2D mode} 
 & \multicolumn{3}{c}{Secondary instability} 
 & \multicolumn{3}{c}{Imminent breakdown} 
 & \multicolumn{3}{c}{End of transition} 
 & \multicolumn{3}{c}{Turbulent} \\
\cmidrule(lr){2-4} \cmidrule(lr){5-7} \cmidrule(lr){8-10} \cmidrule(lr){11-13} \cmidrule(lr){14-16}
 & CPG & TEQ & TNE & CPG & TEQ & TNE & CPG & TEQ & TNE & CPG & TEQ & TNE & CPG & TEQ & TNE \\
\midrule
$x/\delta^*_0$ & 256 & 252 & 292 & 312 & 291 & 361 & 400 & 337 & 436 & 505 & 388 & 505 & 789 & 600 & 600 \\
$Re_x \times 10^{-6}$ & 4.16 & 4.29 & 4.53 & 4.72 & 4.68 & 5.22 & 5.60 & 5.14 & 5.97 & 6.65 & 5.65 & 6.66 & 9.49 & 7.77 & 7.61 \\
$Re_\tau$ & 138 & 145 & 189 & 160 & 151 & 166 & 387 & 175 & 240 & 514 & 254 & 514 & 913 & 687 & 824 \\
$Re^*_\delta \times 10^{-4}$ & 1.61 & 1.56 & 1.68 & 1.72 & 1.63 & 1.80 & 1.87 & 1.71 & 1.93 & 2.04 & 1.79 & 2.04 & 2.43 & 2.10 & 2.18 \\
$Re_\theta$ & 1177 & 1132 & 1211 & 1365 & 1200 & 1333 & 1815 & 1294 & 1537 & 2469 & 1437 & 1809 & 4974 & 2123 & 2323 \\
$H$ & 13.62 & 13.22 & 13.96 & 12.20 & 13.11 & 14.00 & 10.64 & 12.59 & 12.69 & 10.33 & 11.73 & 11.03 & 10.24 & 10.04 & 10.71 \\
$M_{a,t,max}$ & 0.10 & 0.08 & 0.09 & 0.11 & 0.08 & 0.07 & 0.16 & 0.08 & 0.09 & 0.15 & 0.09 & 0.14 & 0.14 & 0.11 & 0.15 \\
$-B_q^{TR} \times 10^{-2}$ & 7.17 & 2.93 & 5.11 & 6.23 & 2.88 & 3.78 & 8.98 & 3.00 & 4.72 & 8.24 & 3.41 & 6.99 & 7.37 & 4.21 & 7.23 \\
$B_q^{V} \times 10^{-3}$ & 0.00 & $-5.87$ & 5.58 & 0.00 & $-5.77$ & 5.79 & 0.00 & $-6.01$ & 5.64 & 0.00 & $-6.83$ & 5.69 & 0.00 & $-8.45$ & 5.43 \\
$\tilde{T}_{max}/T_e$ & 3.70 & 3.36 & 3.73 & 3.70 & 3.37 & 3.67 & 3.69 & 3.40 & 3.75 & 3.59 & 3.48 & 3.81 & 3.56 & 3.51 & 3.72 \\
$\delta_{M_t}/\delta^{*}_{0}$ & 0.88 & 0.90 & 0.90 & 0.91 & 0.90 & 0.90 & 0.47 & 0.89 & 0.86 & 0.37 & 0.87 & 0.67 & 0.48 & 0.66 & 0.33 \\
\midrule
$\Delta x^+$ & 7.72 & 5.49 & 6.56 & 7.99 & 5.47 & 5.31 & 12.14 & 5.88 & 6.72 & 11.47 & 6.67 & 10.05 & 10.66 & 8.35 & 10.97 \\
$\Delta y_w^+$ & 0.68 & 0.48 & 0.58 & 0.70 & 0.48 & 0.47 & 1.07 & 0.52 & 0.59 & 1.01 & 0.59 & 0.88 & 0.94 & 0.73 & 0.96 \\
$\Delta y^+_{99}$ & 1.54 & 1.52 & 1.97 & 1.74 & 1.58 & 1.72 & 4.00 & 1.81 & 2.46 & 5.22 & 2.59 & 5.20 & 9.15 & 6.89 & 8.27 \\
$\Delta z^+$ & 3.59 & 2.55 & 3.05 & 3.71 & 2.54 & 2.47 & 5.63 & 2.73 & 3.12 & 5.33 & 3.10 & 4.67 & 4.95 & 3.88 & 5.09 \\
$L_z^+$ & 1004 & 714 & 853 & 1038 & 711 & 690 & 1578 & 764 & 873 & 1491 & 867 & 1306 & 1385 & 1085 & 1426 \\
\bottomrule
\end{tabular}
\label{table:global_quantities_models}
\end{sidewaystable}
\newpage

% \newpage 

\begin{figure}[H]
    \centering
    \includegraphics[width=22.3cm,angle=-90]{resources/figures/breakdown/result_contour_dns_skinfriction_new.pdf}
    \caption[Time-averaged wall skin friction coefficient ($C_f$) during boundary layer transition for (a) CPG, (b) TEQ, and (c) TNE models.]%
    {Time-averaged wall skin friction coefficient ($C_f$) during boundary layer transition for (a) \gls{CPG}, (b) \gls{TEQ}, and (c) \gls{TNE} models.}
    \label{figure:thermal_spanwise_st}
\end{figure}

\newpage


Further evidence of three-dimensionality is given by the spanwise Stanton-number distributions in Figure~\ref{figure:thermal_spanwise_st}. At early stages, the wall heat flux remains nearly uniform across the span, consistent with the dominance of two-dimensional second-mode disturbances. As transition progresses, streak-like structures emerge, redistributing heat flux into alternating high- and low-intensity bands. These patterns mark the growth of secondary instabilities that destabilise the initially two-dimensional Mack modes and drive rapid breakdown. By the end of the transitional region, the streak organisation weakens and the wall heat flux becomes more evenly distributed, consistent with a fully developed turbulent state.



\subsubsection{Transition structures}
The initial phase of transition is characterised by the amplification of two-dimensional disturbances, dominated by the Mack second mode. Evidence of this process is seen in the small bulge in the skin-friction and Stanton number distributions upstream of breakdown, which reflects the resonance of the primary instability wavepacket. Although Fourier analysis of disturbance signals is not presented here, the observed growth trends are consistent with predictions from linear stability theory, where the second mode exhibits the highest amplification. These early signatures therefore represent the onset of modal growth prior to the emergence of secondary instabilities and subsequent breakdown. The following section also requires the inclusion of isosurfaces for the cases to show what kind of breakdown is happening. Dwell more into the transition dynamics. 


\begin{figure}[H]
    \centering
    \includegraphics[width=10.3cm,angle=0]{resources/figures/test/full_isosurfaces.pdf}
    \caption{Linear stability analysis at $\delta_{*}^{0} = 255$ for (a) CPG, (b) TEQ, and (c) TNE models.}
    \label{figure:thermal_spanwise_st}
\end{figure}

\begin{figure}[H]
    \centering
    \includegraphics[width=10.3cm,angle=0]{resources/figures/test/test_plot.png}
    \caption{Linear stability analysis at $\delta_{*}^{0} = 255$ for (a) CPG, (b) TEQ, and (c) TNE models.}
    \label{figure:thermal_spanwise_st}
\end{figure}


\newpage

\begin{figure}[H]
    \centering
    \includegraphics[width=22.3cm,angle=-90]{resources/figures/breakdown/result_contour_yplus1_uvelocity.pdf}
    \caption[Instantaneous ($u/u_e$) during boundary layer transition for (a) CPG, (b) TEQ, and (c) TNE models.]%
    {Instantaneous ($u/u_e$)  during boundary layer transition for (a) \gls{CPG}, (b) \gls{TEQ}, and (c) \gls{TNE} models.}
    \label{figure:thermal_spanwise_st}
\end{figure}

\newpage



\subsubsection{Forcing interaction}


\subsection{Turbulent regime}
Once the boundary layer has undergone transition, the flow enters a fully turbulent regime characterised by intensified momentum and energy transport. In this stage, wall quantities such as skin friction and heat flux settle towards their asymptotic turbulent levels, while differences between thermal models primarily influence the turbulence structure and mean-flow development. To examine these effects, velocity statistics and Reynolds stresses are investigated. Due to the relatively small range of Reynolds numbers attained in the turbulent portion of the domain, the nondimensional flow profiles do not vary substantially along the flat plate. Therefore, the analysis is restricted to a single representative station. The corresponding streamwise locations for each model are listed in Table~\ref{table:global_quantities_models} in the turbulent region column.



\subsubsection{Velocity and temperature statistics}
A detailed assessment of velocity statistics is essential for characterising the turbulent boundary layer. Mean velocity profiles provide direct information on the redistribution of momentum by turbulence and form the basis for testing scaling laws and transformation approaches. Higher-order statistics, such as Reynolds stresses, quantify the intensity and anisotropy of turbulent fluctuations and reveal how energy is transferred among flow components. Together, these measures offer a comprehensive view of the turbulent structure and allow differences between thermal models to be identified and interpreted. In compressible boundary layers, direct comparison with incompressible velocity profiles requires a transformation that accounts for density and viscosity variations near the wall. The van Driest II transformation provides such a framework by extending the original van Driest scaling to include viscous effects, offering improved collapse of velocity profiles in the inner region. This makes it a common reference for assessing compressible turbulent boundary layers and for comparing different thermal models on a consistent basis.  The mean velocity and wall-normal coordinate $y$ are scaled using the friction velocity $u_{\tau}$ and the viscous length scale $\nu_w/u_{\tau}$. The van Driest transformation is commonly employed to compare compressible turbulent boundary layers with their incompressible counterparts. The transformed velocity is defined as  
\begin{equation}
    u^{+} = \int_{0}^{\bar{u}} \sqrt{\frac{\bar{\rho}}{\rho_w}} \, du ,
\end{equation}
where $\bar{\rho}$ is the time averaged local mean density and $\rho_w$ the density at the wall.  

The resulting mean velocity profile should collapse onto the universal law of the wall, matching the viscous sublayer ($u^{+}=y^{+}$ for $y^{+}<10$) and approaching the logarithmic law at larger wall-normal distances:
\begin{equation}
    u^{+} = \frac{1}{\kappa}\ln(y^{+}) + C, 
    \qquad \kappa = 0.41, \quad C = 5.2 .
\end{equation}

Recent work has introduced updated velocity transformations that improve the collapse of compressible profiles, particularly under cold-wall conditions~\cite{HasanEtal2023}. These approaches build on the semilocal scaling ideas of Trettel and Larsson~\cite{TrettelLarsson2016}, which emphasise the role of variable density and viscosity in defining an effective viscous length scale.

To define the transformation, a consistent relationship between $Y^{+}$ and $y^{+}$ is required. Under the assumption of universality of the turbulent shear stress in the near-wall region, the relation
\begin{equation}
    Y^{+} = \frac{\delta_{v}}{\delta^{*}_{v}}\, y^{+} = y^{*}
\end{equation}
is introduced. In practice, this universality is not guaranteed once compressibility modifies the near-wall stress distribution, so the equivalence $Y^{+}=y^{*}$ requires reassessment. The interpretation
\begin{equation}
    Y^{+} = \frac{\nu}{\delta^{*}_{v}}
\end{equation}
identifies $\delta^{*}_{v}$ as the relevant viscous length scale for small-scale motions in compressible flows and forms the basis of the HLPP formulation.

Making use of the coordinate mapping $Y^{+}=y^{*}$, together with
\begin{equation}
    \frac{dy^{*}}{dy} = \left(1-y^{*}\frac{d\delta^{*}_{v}}{dy}\right)\frac{\delta_{v}}{\delta^{*}_{v}},
    \qquad
    \frac{u_{\tau}}{u^{*}_{\tau}}=\sqrt{\frac{\bar{\rho}}{\rho_{w}}},
\end{equation}
the transformation kernel becomes
\begin{equation}
\frac{d\overline{U}^{+}}{dy^{+}} = 
\left(\frac{1+\mu^{c}_{t}/\bar{\mu}}{1+\mu^{i}_{t}/\mu_{w}}\right)^{1/3}
\left(1-\frac{y}{\delta^{*}_{v}}\frac{d\delta^{*}_{v}}{dy}\right)^{1/2}
\sqrt{\frac{\bar{\rho}}{\rho_{w}}},
\label{eq:TLkernel}
\end{equation}
where the three multiplicative terms represent adjustments to (i) friction velocity scaling, (ii) the viscous length scale, and (iii) density variations in the inner layer.

To close the formulation, explicit models for the eddy viscosities $\mu^{i}_{t}$ and $\mu^{c}_{t}$ are introduced. The Johnson–King model with Van Driest-type damping is adopted, giving
\begin{equation}
    \mu^{i}_{t}=\sqrt{\tau_{w}\rho_{w}}\,\kappa y D^{i},
    \qquad
    D^{i}=\left[1-\exp\left(-\frac{y^{+}}{A^{+}}\right)\right]^{2},
\end{equation}
with $\kappa=0.41$ and $A^{+}=17$ yielding a log-law intercept close to 5.2. For compressible flows, the corresponding form is
\begin{equation}
    \mu^{c}_{t}=\sqrt{\tau_{w}\rho_{w}}\,\kappa y D^{c},
\end{equation}
where $D^{c}$ is a damping function based on the semilocal coordinate $y^{*}$.

Integrating the kernel gives the transformed velocity profile,
\begin{equation}
u^{+}_{HLPP}(y^{*}) = \int_{0}^{y^{*}}
\left(\frac{1+\mu^{c}_{t}/\bar{\mu}}{1+\mu^{i}_{t}/\mu_{w}}\right)^{1/3}
\left(1-\frac{y}{\delta^{*}_{v}}\frac{d\delta^{*}_{v}}{dy}\right)^{1/2}
\sqrt{\frac{\bar{\rho}}{\rho_{w}}}\, dy^{*},
\label{eq:hlpp}
\end{equation}
where $u^{+}_{HLPP}$ denotes the transformed velocity in wall units. This representation allows mean velocity profiles in compressible boundary layers to be compared consistently across thermal models and remains accurate even under cold-wall conditions.


\begin{figure}[t]
    \centering
    \fontsize{9}{10}\selectfont
    \includegraphics[width=1.04\textwidth]{resources/figures/breakdown/result_thermal_velocity_temperature_statistics.pdf}
    \caption[Velocity and temperature statistics in transformed wall units for Mach~6, $T_w=1800$ K.]%
    {Profiles of velocity and temperature statistics plotted against the semilocal wall coordinate $y^{*}$, obtained from the van Driest transformation, for the Mach~6, $T_w=1800$ K case. Results are shown for \gls{CPG}, \gls{TEQ}, and \gls{TNE} formulations.}
    \label{figure:thermal_statistics}
\end{figure}

As shown in Figure~\ref{figure:thermal_spanwise_st}, the velocity profiles collapse well onto the expected scaling laws when expressed in terms of the semilocal wall coordinate $y^{*}$. The curves follow the linear relation in the viscous sublayer and approach the logarithmic law in the overlap region, confirming that the HLPP transformation provides a consistent framework across the different thermal models.  

The temperature profiles, however, do not show full equilibration. Whereas linear stability theory indicated that in the downstream region the \gls{TNE} mean temperature tends to align more closely with \gls{TEQ} than with \gls{CPG}, the present turbulent results reveal the opposite behaviour: turbulence slows down the equilibration process, and the \gls{TNE} profiles remain distinctly separated from \gls{TEQ} further into the boundary layer. This highlights the role of turbulence in delaying thermal relaxation compared to the laminar or transitional regimes. Overall, these results demonstrate that while velocity statistics conform to classical scaling once turbulence is established, thermal nonequilibrium effects persist much further downstream, underscoring the fundamentally different relaxation behaviour of velocity and temperature fields in turbulent boundary layers.


\begin{figure}[t]
    \centering
    \fontsize{9}{10}\selectfont
    \includegraphics[width=1.04\textwidth]{resources/figures/breakdown/result_thermal_reynolds_stresses.pdf}
    \caption[Reynolds stress components in semilocal wall units for Mach~6, $T_w=1800$~K.]%
    {Reynolds stress components $\overline{\rho}\,\widetilde{u''u''}$, 
     $\overline{\rho}\,\widetilde{v''v''}$, 
     $\overline{\rho}\,\widetilde{w''w''}$, and 
     $-\overline{\rho}\,\widetilde{u''v''}$ plotted against the semilocal wall coordinate $y^{*}$ 
     for the Mach~6, $T_w=1800$~K case. Results are compared between \gls{CPG}, \gls{TEQ}, and \gls{TNE} formulations at multiple streamwise stations, with the final location highlighted.}
    \label{figure:thermal_reynolds_stresses}
\end{figure}

\subsubsection{Reynolds stresses}

The Reynolds stress components at $x/\delta^{*}_{0}=650$ and $700$, extracted from Favre-averaged statistics and normalised by the mean wall shear stress $\tau_{w}$, exhibit the canonical ordering $\overline{\rho}\,\widetilde{u''u''} > \overline{\rho}\,\widetilde{v''v''} \sim \overline{\rho}\,\widetilde{w''w''}$, with the shear stress $-\overline{\rho}\,\widetilde{u''v''}$ peaking closer to the wall, as shown in figure~\ref{figure:thermal_reynolds_stresses}. The streamwise component reaches $\overline{\rho}\,\widetilde{u''u''}/\tau_{w}\approx 14$ at $y^{*}\approx 12$, while $\overline{\rho}\,\widetilde{v''v''}$ and $\overline{\rho}\,\widetilde{w''w''}$ attain peak values of $\mathcal{O}(1.0$–$1.5)$, and the shear stress reaches a maximum of about unity. These observations are consistent with canonical DNS of hypersonic turbulent boundary layers \citep{ZhangDuanChoudhari2018,XuEtal2021}. At both stations, the \gls{CPG}, \gls{TEQ}, and \gls{TNE} closures collapse onto nearly indistinguishable profiles, demonstrating that once turbulence is fully developed, the inclusion of thermal nonequilibrium effects does not significantly alter the self-similar Reynolds stress structure.

\begin{figure}[t]
    \centering
    \fontsize{9}{10}\selectfont
    \includegraphics[width=1.04\textwidth]{resources/figures/breakdown/result_thermal_reynolds_rho_T.pdf}
    \caption[Density and temperature statistics in semilocal wall units for Mach~6, $T_w=1800$~K.]%
    {Profiles of Favre-averaged density (top row) and temperature (bottom row) together with their corresponding normalised root-mean-square fluctuations, plotted against the semilocal wall coordinate $y^{*}$ for the Mach~6, $T_w=1800$~K case. Blue lines correspond to the \gls{CPG} formulation and red lines to the \gls{TNE} formulation, with results shown at several streamwise stations and the final location emphasised.}
    \label{figure:thermal_rho_T}
\end{figure}

Figure~\ref{figure:thermal_rho_T} presents the evolution of density and temperature statistics across the boundary layer. In the mean profiles (left panels), both \gls{CPG} and \gls{TNE} display broadly similar behaviour, with the density and temperature ratios decreasing from elevated near-wall values towards the freestream. Subtle differences emerge downstream, where the \gls{TNE} solutions relax more gradually, reflecting the delayed equilibration associated with vibrational energy modes.

More pronounced deviations are observed in the fluctuation levels (right panels). Within the turbulent region, the \gls{TNE} formulation consistently exhibits larger root-mean-square variations in both density and temperature compared to the \gls{CPG} case. This amplification arises from the additional degrees of freedom introduced by vibrational nonequilibrium, which enhance the coupling between thermal and momentum fluctuations. The resulting increase in density and temperature variances strengthens turbulent mixing and alters the redistribution of energy within the boundary layer. In the mean temperature profiles, the \gls{CPG} and \gls{TNE} cases exhibit similar near-wall behaviour, but the \gls{TNE} solution decays more gradually towards the freestream, reflecting delayed vibrational relaxation. These findings indicate that, even when mean profiles appear to converge, the turbulence structure remains sensitive to the inclusion of vibrational nonequilibrium, with direct consequences for predicting wall heat transfer and informing model development.






% \subsection{Transitional regime}



% \section{Role of Mach number and wall thermal conditions}

% \subsubsection{Results}
% \subsection{Breakdown to turbulence}
% Turbulent mean and flow fluctuations
% \subsubsection{Thermal nonequilibrium effects}
% \subsubsection{Heat flux and Reynolds stresses}

% \section{Role of Mach number and wall thermal conditions}

% % \subsection{Results}
% \subsection{Global flow properties}
% \subsection{Boundary-layer profiles}
% % \subsection{Shock--boundary layer interaction}

% \subsection{Breakdown to turbulence}
% \subsubsection{Onset location and growth rates}
% \subsubsection{Mode interactions and energy transfer}
% \subsubsection{Thermal nonequilibrium effects}
% \subsubsection{Heat flux and Reynolds stresses}
% \subsubsection{Second-order statistics}
% % \subsubsection{Spectral analysis of turbulence}

% \newpage

\section{Role of Mach number and wall thermal conditions}
The combined influence of Mach number and wall thermal conditions is now assessed within the framework of thermal nonequilibrium. These parameters exert a strong control on the stability and transition characteristics of hypersonic boundary layers: increasing Mach number enhances compressibility effects and amplifies acoustic-trapped disturbances, while variation in wall temperature modifies the mean thermal gradients that drive instability growth. To isolate these effects, simulations are performed for \gls{TNE} flows at Mach~6 and Mach~8, each with both a cold wall at $T_w=1200$~K and a hot wall at $T_w=1800$~K. This configuration enables a systematic examination of how freestream Mach number and wall thermal state interact to influence instability amplification, breakdown mechanisms, and the resulting surface heating environment.


\subsection{Global flow properties}

\subsection{Breakdown to turbulence}
\subsection{Onset location and growth rates}
% \subsubsection{Boundary-layer profiles}

\subsection{Mode interactions and energy transfer}
\subsection{Streak breakdown}

\subsection{Thermal nonequilibrium effects}
\subsubsection{Heat flux and Reynolds stresses}
\subsubsection{Second-order statistics}



